%!TEX root = ../main.tex

\chapter*{Abstrakt}

Geschäftskritische, verteilte Legacy-Systeme werden häufig über ihren ur\-sprüng\-lich geplanten Lebenszyklus hinaus betrieben und weisen dabei strukturelle Design Smells auf, die Skalierbarkeit und Wartbarkeit beeinträchtigen. Direkte Eingriffe in den Quellcode solcher Systeme sind aufgrund fehlender Dokumentation und mangelnder Expertise häufig mit hohem Risiko verbunden. Eine vorangegangene Arbeit adressierte dieses Problem für ein reales Alarmempfangssystem mittels des \gls{lps}, das synchrone Kommunikationsketten durch eine transparente, \gls{mqtt}-basierte Proxy-Broker-Proxy-Architektur asynchron entkoppelt, ohne den bestehenden Quellcode zu verändern. Eine experimentelle Analyse im Rahmen dieser Masterarbeit zeigt jedoch, dass das \gls{lps} unter realistischen Fehlerbedingungen (Netzwerkfehler, Serverfehler) zentrale Delivery-Garantien verletzt: In den durchgeführten Experimenten traten bei $5\,000$ Alarmmeldungen $77{,}8\,\%$ Duplikate und $13{,}5\,\%$ Reihenfolgeverletzungen auf. Zusätzlich zeigt sich, dass die Skalierbarkeit des \gls{lps} primär durch das Hinzufügen weiterer Service-Instanzen erreicht wird, was den operativen Aufwand erhöht.

Aufbauend auf diesen Erkenntnissen entwickelt diese Masterarbeit das \gls{alps}, eine erweiterte Architektur, die auf \gls{amqp}-basierten Broker-Mechanismen wie Consistent-Hash-Routing, \texttt{x-single-active-consumer} und strikter Flusskontrolle (\texttt{prefetch\_count = 1}) beruht, um eine deterministische \textit{In-Order-Delivery} pro Alarmanlage sicherzustellen. Ergänzt wird die Architektur um eine zu\-sätz\-li\-che, Mnesia-gestützte Deduplizierungskomponente, die als vorgelagerte Sidecar-Komponente fachliche Duplikate erkennt und damit \textit{Exactly-Once-Delivery} Ende-zu-Ende gewährleistet. Die experimentelle Evaluation unter identischen Fehlerbedingungen zeigt, dass \gls{alps} beide Delivery-Garantien vollständig einhält ($0{,}0\,\%$ Duplikate, $0{,}0\,\%$ Reihenfolgeverletzungen) und zu\-sätz\-lich durch Queue-Partitionierung eine effizientere, operativ weniger aufwändige Skalierung ermöglicht als das ursprüngliche \gls{lps}. Die Ergebnisse belegen, dass \gls{alps} eine robuste und skalierbare Lösung für das Architektur-Reengineering geschäftskritischer, verteilter Legacy-Systeme darstellt, ohne Änderungen am bestehenden Quellcode zu erfordern.

\chapter{Einleitung}
\label{ch:einleitung}

Viele geschäftskritische Legacy-Systeme werden weit über ihren ursprünglich geplanten Lebenszyklus hinaus betrieben und stetig weiterentwickelt. % Review-Punkt 3: „der Technik“
Obwohl die zugrunde liegenden Programmstrukturen und der Stand der Technik oft nicht mehr aktuell sind, bilden diese Systeme weiterhin ein unverzichtbares Rückgrat in produktiven Umgebungen~\citep{363157}. Solche Systeme weisen häufig Architekturmerkmale auf, die Skalierbarkeit, Wartbarkeit und Performance negativ beeinflussen. % Review-Punkt 2: Formulierung glätten (Weiterentwicklung; aus historischen Gründen)
Historisch bedingt konzentriert sich ihre Weiterentwicklung meist auf die Erfüllung funktionaler Anforderungen und ließ zukünftige Herausforderungen bezüglich Skalierbarkeit und Wartbarkeit außer Acht~\citep{refactoring_for_software_architecture_smells}. Dieser Prozess begünstigt die Entstehung sogenannter Design Smells.

Im Gegensatz zu direkten Programmfehlern sind diese Smells Indikatoren für tiefere strukturelle Mängel, die die Verständlichkeit, Wartbarkeit und Testbarkeit eines Systems beeinträchtigen~\citep{software_design_smell_detection}. Sie entstehen oft durch bewusste Abkürzungen zugunsten einer schnelleren Auslieferung. Design Smells wirken sich negativ auf die weitere Entwicklung aus, verschlechtern die Wartbarkeit und erschweren die Erfüllung wachsender funktionaler Anforderungen~\citep{sharma2020}.

Die Bewältigung dieser angesammelten technischen Herausforderungen erfordert Reengineering und Refactoring, um die Codebasis wieder an aktuelle funktionale sowie nicht-funktionale Qualitätsansprüche anzupassen~\citep{laguna_systematic_2013, majthoub_software_2018}.

Ein Reengineering kann tiefgreifende Änderungen am Quellcode umfassen oder sich auf die Neustrukturierung der zugrunde liegenden Architektur konzentrieren~\citep{laguna_systematic_2013}. % Review-Punkt 15: einheitlich „klassisch“ statt „traditionell“
% Review-Punkt 4: kein „traditionell“ im Refactoring-Kontext
% Review-Punkt 5: Abgrenzung Reengineering vs. Code-Refactoring; Risiko nicht nur beim Code-Refactoring
Während klassisches \textit{Code-Refactoring} darauf abzielt, die interne Codestruktur ohne Änderung des beobachtbaren Verhaltens zu verbessern, erweisen sich sowohl tiefgreifende Code-Eingriffe als auch architektonische Reengineering-Maßnahmen im industriellen Umfeld häufig als zeitaufwändig und riskant~\citep{stevens_systems_1998, matthiesen_why_2015}. Code-Refactoring bleibt dabei lokal auf den Codebestand bezogen; Reengineering kann darüber hinaus Modelle und Architektur umfassen.

Diese Masterarbeit betrachtet eine spezifische Klasse verteilter Legacy-Sys\-te\-me mit den folgenden Eigenschaften: Es handelt sich um langlebige, geschäfts\-kritische Systeme im Produktivbetrieb mit einer umfangreichen, jedoch kaum dokumentierten Codebasis. Aufgrund der fehlenden Dokumentation wird das interne Systemverhalten nicht mehr im Detail verstanden. Zudem steht die Expertise der ursprünglichen Entwickler nicht mehr zur Verfügung. Dadurch stellt jede Code-Änderung ein erhebliches Risiko dar, weil Änderungen an einer Stelle unvorhersehbare Seiteneffekte in anderen Teilen des Systems auslösen können.

Aus diesem Grund konzentriert sich diese Masterarbeit auf die strukturelle Optimierung der Systemarchitektur, um die Auswirkungen der Design Smells zu reduzieren, ohne den bestehenden Quellcode anzupassen. Ziel ist es, die Systemarchitektur so zu transformieren, dass sie aktuellen Anforderungen an Skalierbarkeit, Wartbarkeit und Antwortzeiten gerecht wird, während die bestehenden funktionalen Anforderungen erhalten bleiben.

\section{Verwandte Arbeiten}
\label{sec:related_work}

In der Arbeit von~\citep{liebig_2025_architectural_refactoring} wurde ein reales, verteiltes Legacy-Alarmempfangssystem der Firma GS electronic Gebr. Schönweitz GmbH~\citep{gselectronic} analysiert. Die Kernaufgabe dieses Systems besteht darin, Alarmmeldungen externer Alarmanlagen zu empfangen, zu verarbeiten und in einer Datenbank zu persistieren. Die Meldungen werden anschließend für Mitarbeiter in einer Leitstelle bereitgestellt. In der bestehenden Architektur durchlaufen die Meldungen eine synchrone, sequentielle Verarbeitungskette. Das System ist dabei auf \textit{Exactly-Once-Delivery} und \textit{In-Order-Delivery} angewiesen. Diese Eigenschaften sowie die zugrunde liegende Funktionsweise werden in Kapitel~\ref{ch:fallstudie} beschrieben.

Im Rahmen dieser Analyse identifizierten die Autoren fünf spezifische Design Smells, die Architektur und Verhalten des betrachteten Systems beeinträchtigen. Um ihre Auswirkungen zu reduzieren, führten sie ein Reengineering der Systemarchitektur des Legacy-Systems durch.

% Frühe Einführung LPS
Der Kern dieses Ansatzes besteht in der transparenten Einführung des \gls{lps} zwischen den bestehenden Komponenten des verteilten Legacy-Systems, ohne den Quellcode anzupassen. Durch den Übergang von sequentieller zu asynchroner Verarbeitung innerhalb des \gls{lps} konnten die Antwortzeiten reduziert und die identifizierten Design Smells adressiert werden. Die bestehenden funktionalen Systemanforderungen wurden dabei beibehalten.

Die Evaluation zeigte zudem, dass durch diesen Ansatz die Antwortzeiten des Legacy-Systems verbessert wurden, sodass das System den Vorgaben der europäischen Norm EN 50136~\citep{din_en_50136-12012-08_alarm_2012} entspricht.

Damit stellt der Ansatz einen wichtigen Fortschritt dar, löst die Problemstellung jedoch nicht vollständig. Die Einhaltung der Delivery-Garantien wurde in der zugrunde liegenden Arbeit vor allem für den Normalbetrieb gezeigt. Eine umfassende Absicherung unter kritischen Fehlerszenarien sowie eine belastbare Bewertung von Korrektheit und Skalierbarkeit unter Fehlerbedingungen bleiben als offene Punkte bestehen~\citep{liebig_2025_architectural_refactoring}. Genau an diesen offenen Punkten setzt die vorliegende Masterarbeit an.

%Darüber hinaus konnte belegt werden, dass die implementierte Lösung die ursprünglichen Delivery-Garantien des bestehenden Systems vollständig einhält.

\section{Wissenschaftlicher Beitrag}
\label{sec:wissenschaftlicher_beitrag}

% Frühe Einführung ALPS
Aufbauend auf dem \gls{lps}-Ansatz zur strukturellen Optimierung verteilter Legacy-Systeme von~\citep{liebig_2025_architectural_refactoring} entwickelt diese Masterarbeit das \gls{alps} und liefert die folgenden vier Kernbeiträge:

\begin{itemize}
    % Review-Punkt 8: Schwachstellenanalyse vermeiden (Security-Konnotation)
    \item \textbf{Analyse der Limitationen:} Experimentelle Evaluierung des in~\citep{liebig_2025_architectural_refactoring} entwickelten \gls{lps}. Dabei werden kritische Defizite bei der Skalierbarkeit sowie Verletzungen zentraler Delivery-Garantien (\textit{Exactly-Once-Delivery} und \textit{In-Order-Delivery}) in Fehlerszenarien aufgezeigt.

    \item \textbf{Optimierte Architektur:} Konzeption und Implementierung von \gls{alps}, das die Vorteile der asynchronen Verarbeitung bewahrt, eine strikte \textit{In-Order-Delivery} garantiert und eine bessere Skalierbarkeit aufweist.

    \item \textbf{Deduplizierungsmechanismus:} Erweiterung von \gls{alps} um eine Komponente zur Duplikaterkennung mittels einer verteilten Datenbank, um die \textit{Exactly-Once-Delivery} sicherzustellen.

    \item \textbf{Experimentelle Evaluation:} Durchführung und Auswertung der in dieser Masterarbeit definierten Benchmarks und Simulationen für \gls{lps} und \gls{alps}. Die Ergebnisse zeigen Verbesserungen hinsichtlich Skalierbarkeit und Korrektheit unter Beibehaltung der funktionalen Anforderungen.
\end{itemize}

\section{Aufbau der Arbeit}

% Review-Punkt 7: Kapitelbeschreibung inkl. Implementierung von ALPS
Die Masterarbeit ist wie folgt strukturiert: Kapitel~\ref{ch:einleitung} motiviert die Problemstellung und ordnet den wissenschaftlichen Beitrag ein. Kapitel~\ref{ch:grundlagen} vermittelt die theoretischen Grundlagen, die für das weitere Verständnis dieser Arbeit erforderlich sind. Kapitel~\ref{ch:fallstudie} beschreibt das betrachtete Legacy-System, die identifizierten Design Smells sowie das originale \gls{lps}. Kapitel~\ref{ch:amqp-proxy} stellt Design und Implementierung des entwickelten \gls{alps} vor. Kapitel~\ref{ch:eval} enthält die experimentelle Evaluation, und Kapitel~\ref{ch:vergleich} diskutiert die Ergebnisse im Vergleich zum originalen \gls{lps}.

\chapter{Grundlagen}
\label{ch:grundlagen}

\section{Design Smells}
\label{sec:grundlagen_design_smells}

Design Smells bezeichnen strukturelle Auffälligkeiten in einem Softwaresystem, die nicht zwingend zu unmittelbaren Laufzeit- oder Übersetzungsfehlern führen, aber Qualitätsattribute wie Verständlichkeit, Testbarkeit, Erweiterbarkeit und Wartbarkeit negativ beeinflussen können~\citep{software_design_smell_detection}. Sie dienen damit als Hinweise auf Wartbarkeitsprobleme und mögliche Ansatzpunkte für Refactoring oder Reengineering.

Solche Smells können auf unterschiedlichen Ebenen auftreten. Neben lokalen Implementierungsproblemen gibt es Design- und Architektur-Smells, die mehrere Komponenten betreffen und dadurch Auswirkungen auf das Gesamtsystem haben können~\citep{refactoring_for_software_architecture_smells}. Für verteilte Legacy-Systeme sind insbesondere Smells relevant, die aus Kommunikationsstrukturen und Redundanzentscheidungen entstehen, etwa synchrone Request-Response-Ketten~\citep{tiwari_unveiling_2024, liebig_2025_architectural_refactoring} oder heterogene Redundanzmodi~\citep{choudhary_rajesh_high_2025, liebig_2025_architectural_refactoring}.

Da solche Smells häufig nicht isoliert in einzelnen Codeabschnitten liegen, sondern sich in Kommunikationspfaden, Abhängigkeiten und Betriebsmodi mehrerer Services manifestieren, reicht lokales Code-Refactoring nicht immer aus. Ihre Behandlung kann stattdessen ein Reengineering auf Architektur- oder Systemebene erfordern. Der folgende Abschnitt grenzt deshalb Reengineering und Refactoring voneinander ab.

\section{Reengineering und Refactoring}
\label{sec:reengineering_refactoring}

% Review-Punkt 12: Refactoring-Begriffsumfang (Code vs. Architektur) explizit machen
Reengineering und Refactoring sind verwandte Konzepte der Software-Evo\-lu\-ti\-on, unterscheiden sich jedoch im Umfang der vorgenommenen Änderungen. Im engeren Sinn bezeichnet \textit{Code-Refactoring} eine Umstrukturierung eines bestehenden Codebestands, bei der die interne Struktur verbessert wird, ohne das externe Verhalten des Systems zu verändern~\citep{laguna_systematic_2013}. Der Fokus liegt damit auf internen Qua\-li\-täts\-merk\-ma\-len wie Verständlichkeit, Wartbarkeit und Änderbarkeit. Daneben wird \textit{Architectural Refactoring} als architektonische Umstrukturierung verstanden, die Kommunikationspfade, Komponentengrenzen oder Betriebsmodi betrifft und ebenfalls ohne Änderung der fachlich beobachtbaren Funktionalität auskommen kann.

Reengineering ist weiter gefasst. Es umfasst die Untersuchung und Ver\-än\-de\-rung eines bestehenden Systems, um es in eine neue Form zu überführen~\citep{laguna_systematic_2013}. Dies kann Quellcode, Modelle oder Architektur betreffen. Für Legacy-Systeme ist Reengineering besonders relevant, wenn ein System weiterhin geschäftlichen Wert besitzt, aber Änderungen nur schwer oder mit hohem Risiko möglich sind~\citep{stevens_systems_1998}.

% Review-Punkt 13: Mentor-Formulierung zu Architektur-Reengineering
Für diese Masterarbeit ist vor allem Architektur-Reengineering relevant. Die betrachtete Klasse von verteilten Legacy-Systemen soll nicht durch lokale Code-Än\-de\-run\-gen verbessert werden, sondern durch eine strukturelle Anpassung ihrer Architektur. Auf diese Weise vermeidet Architektur-Reengineering das hohe Risiko direkter Code-Änderungen an Legacy-Systemen, indem Eingriffe in die bestehende Implementierung der Fachkomponenten vermieden und stattdessen die Interaktionsstruktur angepasst werden. Dadurch sollen die Auswirkungen der identifizierten Design Smells reduziert werden, während die bestehende Funktionalität erhalten bleibt.

\section{\acrfull{lps}}
\label{sec:grundlagen_lps}

Das \gls{lps} ist ein Architektur-Reengineering-Ansatz für verteilte Legacy-Systeme. Ziel ist es, strukturelle Probleme in der Kommunikation zwischen bestehenden Services zu reduzieren, ohne den Quellcode dieser Services anzupassen~\citep{liebig_2025_architectural_refactoring}. % Review-Punkt 15: „klassisch“ als gewählte Variante (nicht „traditionell“)
Damit unterscheidet sich \gls{lps} von klassischem Code-Refactoring: Nicht die interne Implementierung einzelner Komponenten wird verändert, sondern die Interaktion zwischen ihnen.

% Review-Punkt 16: aktive Formulierung der Proxy-Schicht
Hierzu fügt \gls{lps} eine zusätzliche Proxy-Schicht zwischen bestehenden Service-Komponenten des Legacy-Systems ein. Diese Zwischenschicht entkoppelt synchrone Kommunikation und überführt sie über einen Message Broker in eine asynchrone Verarbeitung~\citep{liebig_2025_architectural_refactoring}. Auf diese Weise können insbesondere Design Smells adressiert werden, die aus synchronen Request-Response-Ketten, Redundanzentscheidungen oder fehlender Entkopplung entstehen.

Die konkrete Architektur, die einzelnen Komponenten des \gls{lps} sowie die Art, wie dadurch die Design Smells des betrachteten Legacy-Systems adressiert werden, werden in Abschnitt~\ref{sec:lps_loesung} beschrieben.

% \section{Asynchrone Architektur}

% In diesem Kontext bedeutet Asynchronität, dass eingehende Meldungen zunächst entgegengenommen und zwischengespeichert werden, bevor ihre Weiterverarbeitung erfolgt (\textit{Store-and-Forward}\footnote{Store-and-Forward bezeichnet ein Verfahren, bei dem Nachrichten zu\-nächst gepuffert und erst danach zeitversetzt weitergeleitet werden.}). Dadurch wird die direkte Abhängigkeit zwischen den beteiligten Services reduziert, was die Auswirkungen synchroner Kopplung in der Verarbeitungskette abschwächen kann~\citep{hohpe2004enterprise, newman2021building}.

\newpage

\section{Messaging-Protokolle in verteilten Architekturen}
\label{sec:grundlagen_messaging}

Messaging-Protokolle bilden das fundamentale Rückgrat sowohl des bestehenden \gls{lps} als auch des in dieser Masterarbeit entwickelten \gls{alps}. Sie dienen der räumlichen und zeitlichen Entkopplung verteilter Systemkomponenten. Anstatt Dienste über synchrone, blockierende Aufrufe (wie \gls{http}) direkt miteinander zu verknüpfen, wird die Kommunikation über einen dedizierten Message Broker geleitet~\citep{hohpe2004enterprise}.


\subsection{\acrfull{mqtt}}
\label{subsec:grundlagen_mqtt}

\gls{mqtt} ist ein leichtgewichtiges Publish/Subscribe-Protokoll, das ursprünglich für Umgebungen mit eingeschränkter Bandbreite und instabilen Netzwerken konzipiert wurde~\citep{thangavel2014performance}. % Review-Punkt 18: Bezug auf originales LPS
% Review-Punkt 19: IoT-Standard-Formulierung entschärft / MQTT als Messaging-Protokoll
Die im originalen \gls{lps} verwendete Protokollversion \gls{mqtt}~5.0 stellt eine umfassende Weiterentwicklung früherer \gls{mqtt}-Versionen dar und führt weitere Features wie verbessertes Session-Management und Fehlercodes ein~\citep{oasis2019mqtt, hivemq_mqtt5_essentials}.

Die Kommunikation folgt einer strikten Client-Broker-Architektur\footnote{Clients kommunizieren nie direkt miteinander, sondern ausschließlich über den Broker, was sie räumlich voneinander entkoppelt~\citep{eugster2003many}.}~\citep{oasis2019mqtt}. Das logische Routing basiert auf einer hierarchischen Topic-Struktur, bei der Subscriber explizit Topics abonnieren und Publisher Nachrichten an diese Topics senden~\citep{oasis2019mqtt, eugster2003many}. Ein zentrales Merkmal von \gls{mqtt} sind die integrierten \gls{qos}-Stufen, die das Verbindungsverbindlichkeits- und Zustellverhalten zwischen den Endpunkten definieren~\citep{oasis2019mqtt}. Die Spezifikation unterscheidet drei diskrete Stufen, um unterschiedliche Kompromisse zwischen Protokoll-Overhead und Verlusttoleranz zu adressieren:

\subsubsection{\gls{qos}~0 (\textit{At-Most-Once})}
Diese Stufe repräsentiert eine unbestätigte \textit{Best-Effort}-Zustellung. Die Nachricht wird vom Publisher ohne lokalen Speicherungs- oder Wiederholungsmechanismus abgesendet. Die erfolgreiche Übermittlung hängt vollständig von der Zuverlässigkeit der zugrunde liegenden \gls{tcp}-Verbindung ab. Da kein Kontrollfluss-Handshake stattfindet, weist \gls{qos}~0 den geringsten Protokoll-Overhead und die geringste Latenz auf~\citep{oasis2019mqtt, hivemq_mqtt5_essentials}.
% Review-Punkt 22: aktive Abbildungsbezüge
Abbildung~\ref{fig:mqtt_qos_0} zeigt diesen zustandslosen Nachrichtenfluss, bei dem das Paket emittiert wird, ohne dass der Broker eine rezeptive Rückmeldung an den Client sendet.

\begin{figure}[hbtp]
    \centering
    \includegraphics[width=0.95\textwidth]{./diagram/mqtt_qos_0.png}
    \caption{Nachrichtenfluss unter \acrshort{mqtt} \acrshort{qos}~0 (\textit{At-Most-Once}).}
    \label{fig:mqtt_qos_0}
\end{figure}

\subsubsection{\gls{qos}~1 (\textit{At-Least-Once})}
Hierbei wird die Zustellung der Nachricht an den Subscriber mindestens einmal garantiert. Der Publisher speichert die Nachricht so lange in seiner lokalen Session, bis eine explizite Bestätigung in Form eines \texttt{PUBACK}-Pakets vom Subscriber eintrifft. Bleibt diese Bestätigung innerhalb eines vordefinierten Timeouts aus, wird die Nachricht erneut übermittelt. Aufgrund transienter Netzwerkfehler oder verloren gegangener \texttt{PUBACK}-Pakete kann dies in der Praxis dazu führen, dass der Subscriber dieselbe fachliche Nachricht mehrfach erhält und verarbeiten muss (Duplikation)~\citep{oasis2019mqtt}.
In Abbildung~\ref{fig:mqtt_qos_1} ist neben dem fehlerfreien Normalbetrieb das Verhalten bei einem solchen Paketverlust dargestellt. Bleibt die Bestätigung aus, initiiert der Client nach Ablauf eines Timeouts einen \textit{Retry}\footnote{Retry bezeichnet hier die erneute Übertragung desselben Pakets nach ausbleibender Bestätigung.} mit gesetztem Duplikat-Flag (\texttt{DUP=1}), was die Entstehung von Duplikaten auf Broker-Seite provoziert.

\begin{figure}[hbtp]
    \centering
    \includegraphics[width=0.95\textwidth]{./diagram/mqtt_qos_1.png}
    \caption{Nachrichtenfluss unter \acrshort{mqtt} \acrshort{qos}~1 (\textit{At-Least-Once}) mit möglicher Duplikatbildung bei ausbleibendem \texttt{PUBACK}.}
    \label{fig:mqtt_qos_1}
\end{figure}

\subsubsection{\gls{qos}~2 (\textit{Exactly-Once})}
Die höchste Sicherheitsstufe garantiert eine exakt einmalige Zustellung auf Protokollebene und eliminiert das Risiko anwendungsseitiger Duplikate. Dies wird über einen zweiphasigen, insgesamt vierstufigen Handshake realisiert. Zunächst sendet der Publisher das \texttt{PUBLISH}-Paket, dessen Erhalt vom Broker mit \textit{Publish Received} (\texttt{PUBREC}) bestätigt wird. Erst nach dieser Absicherung signalisiert der Publisher mittels \textit{Publish Release} (\texttt{PUBREL}) die Freigabe zur Auslieferung, was der Broker verarbeitet und abschließend mit \textit{Publish Complete} (\texttt{PUBCOMP}) bestätigt~\citep{oasis2019mqtt}. Durch diesen erweiterten Koordinationsaufwand steigen jedoch sowohl das zu übertragende Datenvolumen als auch die Latenzzeiten signifikant an~\citep{thangavel2014performance}.
% Review-Punkt 22: aktive Abbildungsbezüge
Abbildung~\ref{fig:mqtt_qos_2} zeigt den exakten Kontrollfluss dieses Mechanismus. Das Diagramm verdeutlicht das Zusammenspiel der Bestätigungs- und Freigabeschritte, welches sicherstellt, dass selbst bei einem Netzwerkausfall während der Bestätigungsphase (z.\,B. Verlust des \texttt{PUBREL}-Pakets) der Broker die Nachricht aufgrund der gespeicherten Paket-\gls{id} als bereits verarbeitet identifiziert und somit keine doppelten Daten an die Anwendungsschicht übergibt.

\begin{figure}[hbtp]
    \centering
    \includegraphics[width=0.95\textwidth]{./diagram/mqtt_qos_2.png}
    \caption{Vierstufiger Handshake unter \acrshort{mqtt} \acrshort{qos}~2 (\textit{Exactly-Once}) zur Duplikatvermeidung.}
    \label{fig:mqtt_qos_2}
\end{figure}

\subsection{\acrfull{amqp}}
\label{subsec:grundlagen_amqp}

% Review-Punkt 21: IoT-Kontext bewusst erst hier (Vergleich MQTT/AMQP), nicht als isolierter „IoT-Standard“ in 2.4.1
Während \gls{mqtt} seinen Ursprung im ressourcenbeschränkten \gls{iot}-Umfeld~\citep{thangavel2014performance} hat und dort häufig eingesetzt wird, wurde das \gls{amqp}~\citep{oasis2012amqp10} explizit für den robusten, transaktionssicheren Nachrichtenaustausch in komplexen Enterprise-Architekturen entwickelt~\citep{vinoski2006advanced, ohara2007toward}.

\subsubsection{Das \gls{amqp}-Komponentenmodell}

Im Gegensatz zum direkten Publish/Subscribe-Muster von \gls{mqtt} etabliert \gls{amqp} eine entkoppelte Architektur, die auf vier zentralen funktionalen Einheiten basiert: \textit{Producer}, \textit{Exchange}, \textit{Queue} und \textit{Consumer}. % Review-Punkt 22: aktive Abbildungsbezüge
% Review-Punkt 26: explizit als typisches Beispiel kennzeichnen
Abbildung~\ref{fig:amqp_topology} zeigt den Ablauf eines typischen \gls{amqp}-Komponentenmodells mit diesen Einheiten.

\begin{figure}[hbtp]
    \centering
    \includegraphics[width=0.95\textwidth]{./img/amqp_topo.pdf}
    % Review-Punkt 26
    \caption{Komponentenmodell und Routingpfad einer typischen \acrshort{amqp}-basierten Broker-Topologie.}
    \label{fig:amqp_topology}
\end{figure}

Ein \textit{Producer} emittiert Nachrichten, adressiert diese jedoch niemals direkt an eine Queue. Stattdessen werden die Nachrichten an einen \textit{Exchange} übergeben. Der Exchange fungiert als programmierbare Routing-Instanz innerhalb des Brokers. Er analysiert die Metadaten der eingehenden Nachricht und leitet diese auf Basis von definierten Regeln, den sogenannten \textit{Bindings} (Bindungen), an eine oder mehrere nachgelagerte \textit{Queues} weiter. Die Queues dienen als persistente oder flüchtige Puffer, welche die Nachrichten so lange vorhalten, bis sie von einem oder mehreren registrierten \textit{Consumer} asynchron abgerufen und verarbeitet werden~\citep{oasis2012amqp10, hohpe2004enterprise}.

% Review-Punkt 22: aktive Abbildungsbezüge
% Review-Punkt 23/25/26: Beispielcharakter mit zwei Queues/Consumern
Abbildung~\ref{fig:amqp_sequence} zeigt beispielhaft den protokollarischen Nachrichtenfluss einer typischen \gls{amqp}-Interaktion. Das Diagramm verdeutlicht die voll\-stän\-di\-ge zeitliche und funktionale Entkopplung: Der Producer schließt seinen Kommunikationszyklus mit der Übergabe an den Exchange ab (Schritt 1 \& 2), während die finale Zustellung und die explizite Bestätigung (Schritt 4 \& 5) asynchron und unabhängig auf der Consumer-Seite stattfinden. Zur Veranschaulichung sind zwei Queues mit den zugehörigen Consumern (Consumer~1 an Queue~1, Consumer~2 an Queue~2) dargestellt; dasselbe Muster lässt sich auf Consumer~$1 \dots n$ verallgemeinern.

\begin{figure}[hbtp]
    \centering
    \includegraphics[width=0.95\textwidth]{./diagram/amqp_seq.png}
    % Review-Punkt 23/25/26
    \caption{Beispielhaftes Sequenzdiagramm eines typischen \acrshort{amqp}-Nachrichtenflusses von der Publikation bis zur Bestätigung.}
    \label{fig:amqp_sequence}
\end{figure}

\subsubsection{Routing-Mechanismen und Exchange-Typen}

Das logische Routing zwischen Exchange und Queue basiert auf dem Zusammenspiel von zwei Schlüsseln: Dem \textit{Routing Key}, mit dem der Producer die Nachricht versieht, und dem \textit{Binding Key}, der bei der Verknüpfung einer Queue mit einem Exchange definiert wird. Die Spezifikation unterscheidet standardmäßig drei primäre Exchange-Typen, um unterschiedliche Routing-Topologien abzubilden~\citep{oasis2012amqp10}:

\begin{itemize}
    \item \textbf{Direct Exchange:}~\citep{rabbitmq_exchange_direct} Dieser Typ realisiert ein präzises Punkt-zu-Punkt-Routing. Eine Nachricht wird ausschließlich an die Queues weitergeleitet, deren \textit{Binding Key} exakt bitweise mit dem \textit{Routing Key} der Nachricht übereinstimmt. % Review-Punkt 28: Consumer statt unklarem Worker-Begriff
          Dieses deterministische Verhalten eignet sich ideal für die gezielte Aufgabenverteilung an dedizierte Consumer-Instanzen~\citep{rabbitmq_tut_work_queues, eip_competing_consumers}.

    \item \textbf{Fanout Exchange:}~\citep{rabbitmq_exchange_fanout} Ein Fanout Exchange ignoriert jegliche Routing Keys und arbeitet nach dem Broadcast-Prinzip. Jede eingehende Nachricht wird bedingungslos in alle Queues kopiert, die an diesen Exchange gebunden sind. Dies ermöglicht die einfache Implementierung von One-to-Many-Kommunikationsmustern~\citep{rabbitmq_tut_pubsub}.

    \item \textbf{Topic Exchange:}~\citep{rabbitmq_exchange_topic} Dieser Typ ermöglicht ein flexible, musterbasiertes Routing. Der Routing Key besteht aus einer durch Punkte getrennten Wortliste (z.\,B. \texttt{regions.na.cities.*}). Bindings können Wildcards wie \texttt{*} (ersetzt exakt ein Wort) oder \texttt{\#} (ersetzt null oder mehrere Worte) enthalten, wodurch komplexe, hierarchische Filterstrukturen realisiert werden können.
\end{itemize}

\subsubsection{Zustellungsüberwachung und Flusskontrolle}

% Review-Punkt 62: explizit AMQP-ACK und AMQP-NACK
Zur Absicherung der fehlerfreien Verarbeitung implementiert \gls{amqp} ein bidirektionales Bestätigungsmodell über \textit{Consumer Acknowledgments}. Der Broker entfernt eine Nachricht erst dann endgültig aus der Queue, wenn der Consumer deren erfolgreiche Verarbeitung mit einem \gls{amqp}-\gls{ack} bestätigt. Tritt während der nachgelagerten Verarbeitung ein Fehler auf, oder bricht die Verbindung zum Consumer ab, kann dieser ein \gls{amqp}-\gls{nack} emittieren. Wird dieses \gls{amqp}-\gls{nack} mit dem Parameter \texttt{requeue=true} versehen, % Review-Punkt 29: explizit Datenverlust (Reordering hier nicht primär)
stellt der Broker die Nachricht unmittelbar wieder an der Spitze der Queue bereit, um einen Datenverlust zu verhindern~\citep{oasis2012amqp10}.

Ergänzend hierzu kontrolliert der Parameter \texttt{prefetch\_count} die Flusskontrolle auf Kanalebene. Er definiert das maximale Gleitfenster (\textit{Sliding Window}) unbestätigter Nachrichten, die ein Consumer gleichzeitig vom Broker halten darf. Ein \texttt{prefetch\_count} von 1 erzwingt eine strikte \textit{Stop-and-Wait-Semantik}\footnote{Stop-and-Wait-Semantik bezeichnet ein Kommunikationsverfahren, bei dem der Sender nach dem Versand einer Nachricht wartet und erst nach Erhalt einer expliziten Bestätigung die nächste Nachricht übermittelt.~\citep{tanenbaum2011computer}}: Der Broker liefert die nächste Nachricht erst dann aus, wenn die vorherige vollständig bestätigt wurde. Dies verhindert die Überlastung einzelner Consumer-Instanzen und bildet die funktionale Grundlage für geordnete Verarbeitungssequenzen.

\subsubsection*{Erweiterungen für deterministisches Routing (Message Partitioning)}
% Review-Punkt 30: Formulierung original/originär
Das originale \gls{amqp}-Protokoll garantiert die Erhaltung der Nachrichtenreihenfolge (\textit{In-Order-Delivery}) laut Spezifikation nur für Nachrichten, die über denselben Kanal publiziert, über einen einzigen Pfad in eine einzige Queue geleitet und von einem einzigen Consumer sequenziell verarbeitet werden~\citep{oasis2012amqp10}. Für Systeme, die eine geordnete Verarbeitung pro Entität bei gleichzeitiger horizontaler Skalierung fordern, reicht dies nicht aus.

Dieses Problem wird durch das Konzept der \textit{Partitioned Logs} gelöst, bei dem Nachrichten anhand eines Schlüssels deterministisch auf Partitionen verteilt werden, um die Reihenfolge pro Schlüssel zu garantieren~\citep{kleppmann2017designing}. In der \textit{RabbitMQ}-basierten~\citep{rabbitmq_broker} Implementierung dieser Masterarbeit wird dieses Muster durch spezifische Broker-Erweiterungen umgesetzt, die über den \gls{amqp}-Standard hinausgehen:

\begin{itemize}
    \item \textbf{Consistent-Hash-Routing (\texttt{x-consistent-hash}):} Ein spezialisierter Exchange-Typ, der Nachrichten nicht nach exakten Mustern verteilt, sondern den Hashwert des Routing Keys berechnet und deterministisch auf eine feste Menge von Queues mappt~\citep{rabbitmq_consistent_hash_exchange}. Dadurch wird garantiert, dass alle Nachrichten derselben Entität in derselben Queue landen.

    \item \textbf{Single Active Consumer (\texttt{x-single-active-consumer}):} Eine Queue-Eigenschaft, die sicherstellt, dass unabhängig von der Anzahl registrierter Consumer stets nur ein einzige Consumer aktiv Nachrichten aus dieser Queue bezieht~\citep{cogoluegnes2022rabbitmq}.
\end{itemize}

In Kombination mit einer strikten Flusskontrolle (\texttt{prefetch\_count = 1}) erlauben diese Mechanismen den Aufbau einer Architektur, die parallele Verarbeitung über das Gesamtsystem hinweg ermöglicht, aber für einzelne Entitäten eine strikte \textit{In-Order-Delivery} garantiert. Der genaue Einsatz von diesen Erweiterungen wird in Kapitel~\ref{ch:amqp-proxy} erklärt.

Obwohl die konkrete Umsetzung und Konfiguration in dieser Masterarbeit auf \textit{RabbitMQ}~\citep{rabbitmq_broker} zugeschnitten ist, handelt es sich bei den verwendeten Konzepten nicht um exklusive Funktionen dieses Anbieters. Vergleichbare Erweiterungen existieren in unterschiedlicher Form auch in anderen gängigen Message Brokern, wodurch ein starker \textit{Vendor-Lock-in}\footnote{Vendor-Lock-in bezeichnet die starke technische Abhängigkeit von einem konkreten Anbieter oder Produkt.} vermieden wird. So bietet beispielsweise auch \textit{LavinMQ}~\citep{lavinmq_broker} native Unterstützung für \textit{Consistent-Hash-Routing}~\citep{lavinmq_consistent_hash} und \textit{Single Active Consumer}~\citep{lavinmq_single_active_consumer}. Ebenso stellt \textit{Apache ActiveMQ Artemis}~\citep{artemis_broker} durch das sogenannte \textit{Message Grouping}~\citep{artemis_message_grouping} ein funktionales Äquivalent zum \textit{Consistent-Hash-Routing} bereit und setzt das Prinzip des \textit{Single Active Consumers} mittels \textit{Exclusive Queues}~\citep{artemis_exclusive_queues} um. Diese Beispiele verdeutlichen, dass die in der Architektur eingesetzten Mechanismen zur Wahrung der \textit{In-Order-Delivery} auf etablierten Paradigmen beruhen, die nicht an einen spezifischen Broker gebunden sind.

\section{Erlang und die \acrshort{beam}-Laufzeitumgebung}
\label{sec:grundlagen_erlang}

Für die Realisierung von hochverfügbaren verteilten Systemen spielen die in\-hä\-ren\-ten Eigenschaften der Ausführungsumgebung eine fundamentale Rolle. Die funktionale Programmiersprache \textit{Erlang}~\citep{armstrong2003making} sowie die zugrunde liegende virtuelle Maschine \gls{beam}~\citep{stenman2015erlang, beam_vm} wurden gezielt für den Einsatz in massiv parallelen, verteilten und fehlertoleranten Telekommunikationssystemen entwickelt~\citep{armstrong2003making, virding1996concurrent}. Im Kontext dieser Masterarbeit bildet \textit{Erlang} das Fundament für die zustandsbehaftete Deduplizierungsschicht, da seine Architektur fundamentale Paradigmen verteilter Systeme nativ implementiert.

Die funktionale Natur von Erlang eliminiert Seiteneffekte auf Sprachebene, was die formale Verifikation und die Vorhersagbarkeit von Systemzuständen unter hoher Last erleichtert~\citep{armstrong2007programming}. Die fundamentale Stärke der \gls{beam}-Ar\-chi\-tek\-tur liegt jedoch in der konsequenten Umsetzung des \textit{Actor-Modells}~\citep{agha1986actors}. Im Gegensatz zu traditionellen, \gls{os}-basierten Thread-Modellen, die auf einem gemeinsamen Speicherbereich (\textit{Shared Memory}) operieren und Synchronisationsmechanismen (wie Mutexe oder Semaphore) erfordern, implementiert Erlang eine strikte \textit{Share-Nothing-Architektur}~\citep{armstrong2003making}. Die Realisierung dieser strukturellen Entkopplung innerhalb der \gls{beam}-Laufzeitumgebung stützt sich auf drei fundamentale Eigenschaften:

\begin{itemize}
    \item \textbf{Leichtgewichtige Prozesse:} \gls{beam}-Prozesse sind isolierte Aus\-füh\-rungs\-ein\-hei\-ten, die von der virtuellen Maschine verwaltet werden und einen minimalen Speicher-Overhead (wenige Kilobyte) aufweisen. Dies erlaubt die graduelle Skalierung auf Millionen simultaner Prozesse innerhalb einer einzigen Instanz~\citep{armstrong2003making, armstrong2007programming}.

    \item \textbf{Asynchroner Nachrichtenaustausch:} Die Kommunikation zwischen Prozessen erfolgt ausschließlich über das zustandslose Kopieren von Nachrichten in die mailbox-Strukturen der Zielprozesse~\citep{agha1986actors}. Dies isoliert den Ausfall eines Prozesses vollständig von seiner Umgebung~\citep{armstrong2003making, virding1996concurrent}.

    \item \textbf{Präemptives Scheduling:} Das Scheduling der \gls{beam}-Prozesse basiert auf einem fairen, präemptiven Algorithmus, der strikt verhindert, dass ein einzelner Prozess die Ausführungsumgebung langfristig blockiert. Die Ausführung wird nach einem vordefinierten Zeitfenster (\textit{Time Slice}) unterbrochen, um einen Kontextwechsel zu erzwingen; typischerweise ist dieses Zeitfenster auf etwa 500 \textit{Reductions} (Funktionsaufrufe) begrenzt~\citep{cesarini2009erlang, christakis2013systematic}. Durch diese Begrenzung wird garantiert, dass rechenintensive Operationen einzelner Komponenten niemals die Latenz des gesamten Kommunikationspfads negativ beeinflussen und verlässliche Antwortzeiten im Millisekundenbereich aufrechterhalten werden.
\end{itemize}

Ergänzt wird die \gls{beam}-Laufzeitumgebung durch das \gls{otp}-Framework, welches formale Entwurfsmuster für Fehlertoleranz bereitstellt~\citep{logan2011erlang}. Das dort verankerte \textit{\enquote{Let-it-crash}}-Prinzip besagt, dass Prozesse bei unerwarteten Fehlzuständen unmittelbar terminieren sollten, anstatt defensive Fehlerbehandlungsroutinen zu durchlaufen~\citep{logan2011erlang}.

\subsection{Das verteilte Datenbanksystem Mnesia}
\label{subsec:grundlagen_mnesia}

Aufbauend auf der Erlang \gls{otp}-Laufzeitumgebung integriert \gls{beam} das verteilte Datenbankmanagementsystem \textit{Mnesia}~\citep{mattsson1999mnesia, mnesia}. Mnesia wurde spezifisch für Anwendungen konzipiert, die eine kontinuierliche Verfügbarkeit, extrem geringe Abfragelatenzen und eine replizierte, transaktionale Datenhaltung über mehrere Netzknoten hinweg erfordern~\citep{mattsson1999mnesia}. Im Rahmen des in dieser Masterarbeit entworfenen \gls{alps} fungiert Mnesia als verteilte \textit{Idempotenz-Datenbank} (Deduplizierungs-Datenbank, engl.\ \textit{Deduplication Store})~\citep{helland2012idempotence}.

Um diese verteilte Architektur zu realisieren, kann Mnesia in einem Cluster operieren, der sich aus mehreren vernetzten Erlang-Knoten (Nodes) zusammensetzt. Innerhalb dieses Verbunds können Datenbanktabellen flexibel über verschiedene Knoten hinweg repliziert werden, wodurch eine hohe Fehlertoleranz (Fault Tolerance) und Datenverfügbarkeit im Falle von Netzwerktrennungen oder Hardwareausfällen gewährleistet wird~\citep{mattsson1999mnesia}. Die Entscheidung, wie die Daten auf den einzelnen Knoten physisch abgelegt werden, hat dabei direkte Auswirkungen auf die Latenz und Persistenz des Systems~\citep{mnesia}.

Für die Gewährleistung der Replikationsintegrität und den Schutz vor Datenverlust bietet Mnesia eine konfigurierbare Zustandshaltung auf Tabellenebene, die sich in drei primäre Speicherformen unterteilt:

\begin{enumerate}
    \item \texttt{\textbf{ram\_copies}}: Die Daten werden ausschließlich im flüchtigen Hauptspeicher gehalten. Dies maximiert den Durchsatz, bietet jedoch keine Persistenzgarantie bei einem vollständigen Cluster-Ausfall.

    \item \texttt{\textbf{disc\_copies}}: Die Tabellenstruktur und die Daten werden sowohl im \gls{ram} gehalten als auch synchron auf ein transaktionales Disk-Log geschrieben. Dies sichert die Daten persistent gegen Knotenausfälle ab.

    \item \texttt{\textbf{disc\_only\_copies}}: Die Daten verbleiben ausschließlich auf dem Festspeichermedium, was den \gls{ram}-Bedarf minimiert, jedoch die Zugriffslatenzen signifikant erhöht.
\end{enumerate}

Gemäß dem \gls{cap}-Theorem~\citep{gilbert2002brewer} agiert Mnesia standardmäßig als \gls{ap}-System, das bei Netzwerkpartitionierungen die Verfügbarkeit über die strikte Konsistenz stellt. Da isolierte Knoten weiterhin Schreibvorgänge akzeptieren, führt ein Netzwerkausfall unweigerlich zu divergierenden Datenbeständen (\textit{Split-Brain-Szenario})\footnote{Da Erlang-Knoten nicht zuverlässig zwischen reinen Netzwerkfehlern und dem tatsächlichen Absturz anderer Knoten unterscheiden können, arbeitet in einem solchen Fall jedes isolierte Teilnetz unabhängig auf seiner lokalen Datenkopie weiter~\citep{cesarini2015designing}.}. Tritt diese Situation ein, bietet Mnesia keine Mechanismen zur automatischen Konfliktauflösung. Das System generiert stattdessen lediglich ein \texttt{inconsistent\_database}-Event und überlässt dem Entwickler die Entscheidung, mit welcher Strategie der konsistente Zustand des Clusters wiederhergestellt werden muss~\citep{mnesia}.

Wie diese Mnesia-basierten Replikations- und Deduplizierungsmechanismen konkret in die Gesamtarchitektur eingebettet sind und wie die Recovery von Split-Brain-Szenarien behandelt wird, wird in Kapitel~\ref{ch:amqp-proxy} erklärt.

\chapter{Fallstudie: Verteiltes Legacy-System}
\label{ch:fallstudie}

Als Fallstudie betrachtet diese Masterarbeit das in~\citep{liebig_2025_architectural_refactoring} untersuchte verteilte Legacy-System \textit{\enquote{\gls{marc}}}~\footnote{Das reale \gls{marc} umfasst mehr als 30 Komponenten; in dieser Masterarbeit wird eine für die Analyse reduzierte Sicht auf die relevanten Kommunikations- und Servicepfade verwendet.} der Firma GS electronic Gebr. Schönweitz GmbH~\citep{gselectronic}. \gls{marc} ist ein verteiltes Alarmempfangssystem, das externe Alarmanlagen überwacht sowie Alarmmeldungen empfängt, verarbeitet und in einer Datenbank persistiert, um sie für die anschließende Bearbeitung durch Mitarbeiter in einer Leitstelle bereitzustellen. Abbildung~\ref{fig:marc_simple_architektur} zeigt eine vereinfachte Architektur von \gls{marc} sowie seine zentrale Komponente \textit{\enquote{\gls{ars}}} zur Verarbeitung von Alarmmeldungen. Ist \gls{ars} nicht verfügbar, können keine Alarmmeldungen bearbeitet werden.

\begin{figure}[hbtp]
    \centering
    \includegraphics[width=0.95\textwidth]{./img/marc_simple.pdf}
    \caption{Vereinfachte Gesamtarchitektur des Legacy-Systems \acrshort{marc}.}
    \label{fig:marc_simple_architektur}
\end{figure}

\section{Architektur}
\label{sec:architektur}

Wie im vorherigen Abschnitt beschrieben, bildet die \gls{ars} die zentrale Komponente des \gls{marc}-Systems. Ihre Hauptaufgabe besteht darin, eingehende Alarmmeldungen zu empfangen, zu verarbeiten und in einer Datenbank zu persistieren. Intern setzt sich die \gls{ars} aus mehreren Services zusammen, die im Folgenden als \gls{srv}~1, \gls{srv}~2 und \gls{srv}~3 bezeichnet werden~\ref{fig:ars_architektur}. \gls{srv}~1 stellt dabei die primäre \gls{http}-Schnittstelle bereit, über die Alarmmeldungen von den externen Alarmanlagen entgegengenommen werden.

Nach dem Empfang einer Alarmmeldung durchläuft sie eine synchrone \textit{Request-Response-Kette}~\footnote{Unter einer Request-Response-Kette wird in dieser Masterarbeit eine sequentielle Folge synchroner Request-Response-Interaktionen über mehrere Services verstanden, bei der jeder nachgelagerte Schritt erst nach erfolgreicher Antwort des vorherigen Schritts fortgesetzt wird.} über alle beteiligten Services bis hin zu \gls{srv}~3, welcher die finale Verarbeitung und die Persistierung der Alarmmeldung in einer Datenbank übernimmt. Die Kommunikation innerhalb dieser Kette erfolgt über eine Kombination aus \gls{http}-Requests und einem proprietären \gls{tcp}-Protokoll. Erst nachdem \gls{srv}~3 die Verarbeitung der Alarmmeldung erfolgreich abgeschlossen hat, wird eine Bestätigung über die VerarbeitungsKette an die sendende Alarmanlage zurückgeleitet. Die Antworten der einzelnen Services dienen ausschließlich als Bestätigung des erfolgreichen Empfangs beziehungsweise der erfolgreichen Verarbeitung und enthalten keine zusätzlichen Nutzdaten. Bei negativen Antworten, Netzstörungen oder Timeouts senden die Alarmanlagen die Alarmmeldungen erneut. Abbildung~\ref{fig:ars_architektur} veranschaulicht die interne Verarbeitungskette der \gls{ars}. Eine eingehende Alarmmeldung wird von \gls{srv}~1 angenommen, synchron über \gls{srv}~2 an \gls{srv}~3 weitergereicht und erst nach erfolgreicher Verarbeitung bestätigt.

\begin{figure}[hbtp]
    \centering
    \includegraphics[width=0.95\textwidth]{./img/ars.pdf}
    \caption{Interne Service-Architektur der Alarm Receiving Software (\acrshort{ars}).}
    \label{fig:ars_architektur}
\end{figure}

\gls{marc} ist über drei Rechenzentren repliziert, wobei in jedem Replikat eine eigene Instanz von \gls{ars} betrieben wird. Während die Services \gls{srv}~1 und \gls{srv}~2 in allen \gls{ars}-Instanzen zur Lastverteilung im Aktiv-Aktiv-Modus operieren, ist \gls{srv}~3 im Aktiv-Passiv-Modus ausgelegt. Folglich ist zu jedem Zeitpunkt nur eine Instanz von \gls{srv}~3 aktiv und für die Verarbeitung von Alarmmeldungen zuständig. Abbildung~\ref{fig:marc_architektur} zeigt die replizierte \gls{marc}-Architektur über drei Rechenzentren und wie die verschiedenen Services miteinander interagieren.

\begin{figure}[hbtp]
    \centering
    \includegraphics[width=0.95\textwidth]{./img/marc.pdf}
    \caption{Replizierte \acrshort{marc}-Architektur über drei Rechenzentren mit Aktiv-Aktiv- und Aktiv-Passiv-Betrieb.}
    \label{fig:marc_architektur}
\end{figure}

\section{Design Smells im Legacy-System}
\label{sec:design_smells}

Die im vorherigen Abschnitt beschriebene Architektur von \gls{marc} führt laut~\citep{liebig_2025_architectural_refactoring} zu mehreren strukturellen Problemen. Die Analyse von~\citep{liebig_2025_architectural_refactoring} identifiziert diese strukturellen Problemen und klassifiziert sie in die folgenden Design Smells:

\begin{itemize}
    \item \textbf{Smell~1}: \textit{Synchrone Request-Response-Ketten}. Die sequentielle Kopplung der \gls{ars}-Services; \gls{srv}~1, \gls{srv}~2 und \gls{srv}~3 führt zu einer Akkumulation der Latenzen. Die Gesamtantwortzeit der Verarbeitung einer Alarmmeldung ergibt sich aus der Summe der Latenzen der Einzelkomponenten.

    \item \textbf{Smell~2}: \textit{Heterogene Redundanzmodi}. Der parallele Einsatz von Aktiv-Aktiv für die \gls{ars}-Services; \gls{srv}~1 und \gls{srv}~2 sowie Aktiv-Passiv für \gls{srv}~3 innerhalb derselben Verarbeitungskette verursacht einen Flaschenhals. Dadurch wird der Gesamtdurchsatz des Systems auf die Kapazität einer einzelnen aktiven \gls{srv}~3-Instanz begrenzt.

    \item \textbf{Smell~3}: \textit{Manuelles Ausfall-Management}. Das System kann Ausfälle nicht autonom beheben, sondern erfordert eine manuelle Überwachung und Intervention durch Mitarbeiter. Dieser manuelle Arbeitsschritt dauert oft Minuten, während Sicherheitsnormen wie die \textit{EN 50136}~\citep{din_en_50136-12012-08_alarm_2012} Reaktionen innerhalb weniger Sekunden fordern. Dadurch ist das System im Ernstfall zu langsam und fehleranfällig.

    \item \textbf{Smell~4}: \textit{Abhängigkeit von Exactly-Once-Delivery}. Das System ist auf eine \textit{Exactly-Once-Delivery} angewiesen, da der Verarbeitung von Alarmmeldungen die \textit{Idempotenz}~\footnote{Idempotenz im Kontext verteilter Systeme bedeutet, dass eine Operation (wie das Verarbeiten einer Nachricht) beliebig oft wiederholt werden kann, ohne dass sich der Systemzustand nach der ersten erfolgreichen Anwendung ändert~\citep{helland2012idempotence}.} fehlt. Das bedeutet intuitiv, dass das System nicht nativ mit Duplikaten umgehen kann, die durch wiederholte Übertragungen der Alarmanlagen entstehen. Mitarbeiter in der Leitstelle müssen Duplikate selbstständig erkennen und ignorieren. Das führt in der Praxis zu Fehlverhalten oder zur gleichzeitigen Bearbeitung derselben Alarmmeldung durch mehrere Mitarbeiter in der Leitstelle.

    \item \textbf{Smell~5}: \textit{Strikte Reihenfolgeabhängigkeit (In-Order-Delivery)}. Das System ist auf die korrekte zeitliche Abfolge von Alarmmeldungen angewiesen, da Alarmmeldungen zustandsbehaftet sind. Aufgrund fehlender Mechanismen zur Neuordnung führen Paketverluste oder hohe Latenzen dazu, dass Alarmmeldungen in der falschen Reihenfolge verarbeitet werden. Ein kritisches Beispiel dafür ist der Empfang einer Entwarnung vor dem eigentlichen Alarm: Wird der verspätete Alarm zuletzt verarbeitet, überschreibt er den Status der Entwarnung. Dies führt zu einer fehlerhaften Interpretation und provoziert unnötige Notfalleinsätze.
\end{itemize}

Abbildung~\ref{fig:exactly_once_diagram} zeigt einen beispielhaften Kommunikationsablauf, der zu einer Verletzung der \textit{Exactly-Once-Delivery} führt: Eine Alarmmeldung wird aufgrund eines Kommunikationsfehlers erneut übertragen und dadurch mehrfach verarbeitet.

\begin{figure}[hbtp]
    \centering
    \includegraphics[width=0.95\textwidth]{./diagram/marc_exactly_once_delivery_violation_sequence.png}
    \caption{Sequenzdiagramm eines beispielhaften Kommunikationsablaufs, der zu einer Verletzung der \textit{Exactly-Once-Delivery} im \acrshort{marc}-System führt.}
    \label{fig:exactly_once_diagram}
\end{figure}

Abbildung~\ref{fig:ein_order_diagram} zeigt einen beispielhaften Kommunikationsablauf, der zu einer Verletzung der \textit{In-Order-Delivery} führt: Eine später erzeugte Alarmmeldung wird vor einer früheren Alarmmeldung verarbeitet, da variable Netzwerlatenzen die Reihenfolge verändern können.

\begin{figure}[hbtp]
    \centering
    \includegraphics[width=0.95\textwidth]{./diagram/marc_in_order_delivery_violation_sequence.png}
    \caption{Sequenzdiagramm eines beispielhaften Kommunikationsablaufs, der zu einer Verletzung der \textit{In-Order-Delivery} im \acrshort{marc}-System führt.}
    \label{fig:ein_order_diagram}
\end{figure}

\section{\acrfull{lps}}
\label{sec:lps_loesung}

\subsection{Architektur}
\label{sub-sec:lps_architektur}

Nachdem in vorherigen Abschnitten das \gls{marc}-System beschrieben wurde, betrachtet dieser Abschnitt das in~\citep{liebig_2025_architectural_refactoring} vorgeschlagene \gls{lps}. Auf Basis der zuvor identifizierten Design Smells in Abschnitt~\ref{sec:design_smells} schlägt die Arbeit~\citep{liebig_2025_architectural_refactoring} ein Ar\-chi\-tek\-tur-Reengineering vor, um die Auswirkungen der identifizierten Design Smells zu mitigieren, ohne Modifikationen am Quellcode der Legacy-Services vorzunehmen. Der Kern dieses Ansatzes ist die transparente Integration des \gls{lps} zwischen den Services des \gls{marc}-Systems. Die Grundidee des \gls{lps} besteht in der Entkopplung der synchronen Verarbeitungskette zwischen den Services \gls{srv}~1 und \gls{srv}~2 durch eine \textit{Proxy-Broker-Proxy-Architektur}~\footnote{Eine Proxy-Broker-Proxy-Architektur bezeichnet ein Integrationsmuster, bei dem ein Message Broker von zwei Proxies gekapselt wird. Dies ermöglicht es, synchrone Kommunikationspfade transparent in eine asynchrone Kommunikationspfade zu überführen, ohne die bestehenden Schnittstellen der ursprünglichen Endpunkte anzupassen~\citep{hohpe2004enterprise}.}. Das \gls{lps} wird zwischen \gls{srv}~1 und \gls{srv}~2 positioniert und setzt sich aus drei Komponenten zusammen: dem \gls{lpone}, einem Message Broker und dem \gls{lptwo}.Abbildung~\ref{fig:lps_architektur} zeigt die transparente Einfügung von \gls{lps} zwischen \gls{srv}~1 und \gls{srv}~2. Die zuvor direkte synchrone Kopplung wird durch die Kette \textit{\gls{lpone} $\rightarrow$ Broker $\rightarrow$ \gls{lptwo}} ersetzt, sodass eingehende Alarmmeldungen zunächst entgegengenommen und im Broker gepuffert werden, bevor sie asynchron weiterverarbeitet werden. Die bestehende Legacy-Schnittstelle bleibt dabei nach außen unverändert.

\begin{figure}[htbp]
    \centering
    \includegraphics[width=0.95\textwidth]{./img/lps.pdf}
    \caption{\acrshort{marc} nach Einfügung von \acrshort{lps} zwischen \acrshort{srv}~1 und \acrshort{srv}~2.}
    \label{fig:lps_architektur}
\end{figure}

In der Proof-of-Concept-Implementierung der Arbeit~\citep{liebig_2025_architectural_refactoring} wird das Protokoll \gls{mqtt}~v5 genutzt, um das \gls{lps} zu realisieren, wobei HiveMQ~\citep{hivemq_2025} als Broker zum Einsatz kommt. Im Normalbetrieb nimmt \gls{lpone} die Requests von \gls{srv}~1 entgegen und publiziert sie auf einem entsprechenden shared Topic, die im Folgendem als Main-Topic bezeichnet wird. Sobald der Broker den Eingang bestätigt, sendet \gls{lpone} unmittelbar eine Bestätigung an das vorgelagerte System zurück. Dadurch wird die synchrone Wartezeit der ursprünglichen Verarbeitungskette reduziert. Parallel dazu konsumiert \gls{lptwo} die Nachrichten aus dem Main-Topic und leitet sie asynchron an \gls{srv}~2 weiter. Abbildung~\ref{fig:lps_normalfluss} zeigt einen beispielhaften Verarbeitungsablauf einer Alarmmeldung innerhalb von \gls{lps} im Normalbetrieb. Unter Normalbetrieb wird die fehlerfreie Verarbeitung einer Alarmmeldung verstanden, bei der weder Netzwerkfehler noch negative Antworten der nachgelagerten Services (\gls{srv}~2, \gls{srv}~3) auftreten.

\begin{figure}[htbp]
    \centering
    \includegraphics[width=0.95\textwidth]{./diagram/lps_normal_sequence.png}
    \caption{Sequenzdiagramm eines beispielhaften Verarbeitungsablaufs einer Alarmmeldung innerhalb von \gls{lps} im Normalbetrieb.}
    \label{fig:lps_normalfluss}
\end{figure}

Um die strengen Delivery-Garantien des Legacy-Systems zu wahren, kombiniert \gls{lps} mehrere Mechanismen. Die Kommunikation zwischen \gls{lpone}, Broker und \gls{lptwo} erfolgt über \gls{mqtt} \gls{qos}~2, was eine \textit{Exactly-Once-Delivery} auf diesem Teilpfad garantiert. Um die \textit{In-Order-Delivery} der Alarmmeldungen zu gewährleisten, werden die \gls{lptwo}-Instanzen so konfiguriert, dass Nachrichten innerhalb derselben Instanz anhand von (\texttt{max\_inflight\_messages = 1}) sequentiell verarbeitet werden. Für Fehlerfälle wie Netzwerkfehler oder negative Antworten der nachgelagerten Services, die hinter \gls{lptwo} auftreten, wird ein \textit{Retry}-Mechanismus~\footnote{Retry bezeichnet das erneute Zustellen einer zuvor fehlgeschlagenen Nachricht.} eingeführt. Bei fehlgeschlagener Übertragung einer Alarmmeldung an die nachgelagerten Services, wird diese Alarmmeldung auf ein separates shared Topic, die im Folgendem als Retry-Topic bezeichnet wird, publiziert. Der Retry-Mechanismus wird durch eine Broker-Erweiterung (Plugin) umgesetzt, die Nachrichten aus dem Retry-Topic strikt priorisiert. Sobald sich Nachrichten im Retry-Topic befinden, wird die Auslieferung regulärer Nachrichten aus dem Main-Topic temporär pausiert. Erst nach der erfolgreichen Verarbeitung der Retry-Nachrichten wird der reguläre Nachrichtenfluss fortgesetzt. Dies verhindert, dass neuere Alarmmeldungen eine ältere, zuvor fehlgeschlagene Alarmmeldung überholen. Der Ablauf dieses Mechanismus gestaltet sich wie folgt:

\begin{enumerate}
    \item \gls{lpone} publiziert eine Alarmmeldung in das Main-Topic.
    \item \gls{lptwo} konsumiert die Alarmmeldung aus dem Main-Topic und versucht sie an \gls{srv}~2 weiterzuleiten.
    \item Im Fehlerfall wird die Alarmmeldung in das Retry-Topic publiziert.
    \item Das Broker-Plugin pausiert die Auslieferung von reguläreren Nachrichten aus dem Main-Topic und priorisiert die Auslieferung der Nachrichten aus dem Retry-Topic.
    \item Nach erfolgreicher Verarbeitung aller Nachrichten im Retry-Topic wird die Auslieferung von reguläreren Nachrichten fortgesetzt.
\end{enumerate}

Abbildung~\ref{fig:lps_retry_priorisierung} zeigt einen beispielhaften Verarbeitungsablauf einer Alarmmeldung innerhalb von \gls{lps} im Fehlerfall.

\begin{figure}[htbp]
    \centering
    \includegraphics[width=0.95\textwidth]{./diagram/lps_retry_prio_sequence.png}
    \caption{Sequenzdiagramm eines beispielhaften Verarbeitungsablaufs einer Alarmmeldung innerhalb von \gls{lps} im Fehlerfall.}
    \label{fig:lps_retry_priorisierung}
\end{figure}

\subsection{Auswirkungen auf die Design Smells}
\label{sub-sec:lps_auswirkungen}

Durch die Integration von \gls{lps} werden die Auswirkungen der zuvor identifizierten Design Smells in Abschnitt~\ref{sec:design_smells} mitigiert:

\begin{itemize}
    \item \textbf{Smell~1}: \textit{Synchrone Request-Response-Ketten}. Die direkte synchrone Kopplung zwischen \gls{srv}~1 und \gls{srv}~2 wird aufgrund der asynchronen Verarbeitung der Alarmmeldungen durch die \textit{Proxy-Broker-Proxy-Architektur} aufgebrochen. \gls{lpone} bestätigt den Eingang einer Alarmmeldung an die vorgelagerten Services direkt nach dem Publizieren der Alarmmeldung an das Broker. Dadurch wird die Latenzakkumulation entlang der Verarbeitungskette reduziert.

    \item \textbf{Smell~2}: \textit{Heterogene Redundanzmodi}. die \textit{Proxy-Broker-Proxy-Architektur} und der Einsatz mehrerer \gls{lptwo}-Instanzen machen die Verarbeitung widerstandsfähiger gegenüber Ausfällen nachgelagerter Services oder Netzwerkverbindungen.

          % TODO: TF is this?
    \item \textbf{Smell~3}: \textit{Manuelles Ausfall-Management}. Zwar bleibt die Notwendigkeit manueller Eingriffe bei schwerwiegenden Ausfällen formal bestehen, jedoch werden die unmittelbaren negativen Auswirkungen auf die Antwortzeiten vorgelagerter Systeme reduziert, da eingehende Alarmmeldungen sicher im Broker zwischengespeichert werden.

    \item \textbf{Smell~4 \& Smell~5}: \textit{Exactly-Once-Delivery} \& \textit{In-Order-Delivery} werden über \gls{qos}~2 und das Retry-Mechanismus adressiert.
\end{itemize}

Zur Validierung des Ansatzes wurden abschließend Experimente durchgeführt, bei denen simulierte Alarmanlagen Alarmmeldungen an das senden. Der Fokus der Experimente lag auf der Reduktion der Antwortzeiten sowie der Einhaltung der \textit{Exactly-Once-Delivery}. Die Experimente belegten, dass die Antwortzeiten des Systems nach dem Integrieren von \gls{lps} sowohl im Normalbetrieb als auch in Fehlerszenarien verbessert wurden und das System die Vorgaben der europäischen Norm EN 50136~\citep{din_en_50136-12012-08_alarm_2012} erfüllt. Die Garantie der \textit{Exactly-Once-Delivery} und \textit{In-Order-Delivery} konnten laut den Experimenten ebenso gewahrt werden.

\subsection{Limitationen der Experimente}
\label{sub-sec:lps_limitationen_experimente}

Nach dem Betrachten der durchgeführten Experimente und der Interpretation der Ergebnisse ergeben sich potentiale Verletzungen der Delivery-Garantien. Laut dem Paper wird die Einhaltung von \textit{Exactly-Once-Delivery} bestätigt, in dem die Anzahl gesendeter Alarmmeldungen der simulierten Alarmanlagen gleich der Anzahl angekommener Alarmmeldung. Jedoch ist der Vergleich der Anzahl nicht ausreichend, um die \textit{Exactly-Once-Delivery} zu bestätigen. Eine beispielhafte Sequenz von Alarmmeldungen, die zur Verletzung von \textit{Exactly-Once-Delivery} sieht folgendermaßen aus: Die Alarmanlage \texttt{A} sendet die Alarmmeldungen \texttt{\#1} und \texttt{\#2}. Alarmmeldung~\texttt{\#1} geht im Netz verloren, während Alarmmeldung~\texttt{\#2} dupliziert ankommt. Die Anzahl gesendeter Alarmmeldungen ist gleich der Anzahl angekommener Alarmmeldungen, jedoch ist Alarmmeldung~\texttt{\#2} mehrfach angekommen. Außerdem wird die globale \textit{In-Order-Delivery} in den Experimenten nicht betrachtet. Innerhalb der \gls{lptwo} Instanzen werden Alarmmeldungen anhand von (\texttt{max\_inflight\_messages = 1}) sequentiell bearbeitet. Dadurch garantiert \gls{lptwo} \textit{In-Order-Delivery} innerhalb einer Instanz, jedoch keine globale \textit{In-Order-Delivery}. Aufgrund der \textit{shared subscriptions} ist es möglich, dass zwei \gls{lptwo} Instanzen Alarmmeldungen gleichzeitig bearbeiten. Das kann bei variablen Verarbeitungszeiten der Alarmmeldungen in den nachgelagerten Services dazu führen, dass Alarmmeldungen in der falschen Reihenfolge in der Datenbank gespeichert werden. Im nächsten Abschnitt werden Experimente durchgeführt, bei denen die Verletzung der Delivery-Garantien gezielt provoziert wird.

\subsection{Experimente}
\label{sub-sec:lps_experimente}

Die in diesem Abschnitt durchgeführten Experimente umfassen zwei Simulationen des Legacy-Systems nach dem Integrieren von \gls{lps}. Bei den Simulationen werden Alarmanlagen simuliert, die Alarmmeldungen periodisch an das System schicken. Die gesendeten Alarmmeldungen bestehen hauptsächlich aus zwei Feldern; \texttt{device\_id} und \texttt{triggertime}. \texttt{device\_id} ist der Identifikator einer Alarmanlage und \texttt{triggertime} ist der Identifikator einer Alarmmeldung. Es ist garantiert, dass die \texttt{triggertime} der einzelnen gesendeten Alarmmeldungen monoton steigend ist. Außerdem sind die simulierten Alarmanlagen sind so eingestellt, dass sie das Übertragen von Alarmmeldungen erneut versuchen, falls Fehler beim Übertragen auftreten. Wiederholte Alarmmeldungen besitzen die gleiche \texttt{triggertime} wie die ursprünglich fehlgeschlagene Alarmmeldung. Nachgelagerte Services von \gls{lps} nämlich \gls{srv}~1 und \gls{srv}~2 werden auch durch einen Service simuliert, deren Aufgabe ist, Alarmmeldungen von \gls{lps} bzw.\,w.\ \gls{lptwo} zu empfangen und speichern. Die gespeicherten Alarmmeldungen werden abschließend auf \textit{Exactly-Once-Delivery} und \textit{In-Order-Delivery} geprüft. Zum Ausführen der Experimente werden vier Hilfsservices eingeführt:

\begin{itemize}
    \item \textbf{Alarmanlagen-Simulator:} Der Alarmanlagen-Simulator bildet externe Alarmanlagen nach und erzeugt periodisch Alarmmeldungen, die an das System geschickt werden. Die gesamte Anzahl simulierter Alarmanlagen, die maximale Anzahl gesendeter Alarmmeldungen pro Alarmanlage sowie die willkürliche Wartezeit zwischen zwei gesendeter Alarmmeldungen sind einstellbar. Erfolgreich gesendeter Alarmmeldungen werden abschließend in einer Datei gespeichert (\texttt{device-results.json}).

    \item \textbf{\gls{tcp}-Proxy:} Der \gls{tcp}-Proxy dient zur gezielten Simulation von Netzwerkfehlern zwischen Komponenten des \gls{lps}, konkret zwischen \gls{lptwo} Instanzen und dem Broker. Der \gls{tcp}-Proxy akzeptiert Verbindungen und leitet sie and die Zieladresse weiter. Die willkürliche maximale Dauer einer akzeptierten Verbindung ist einstellbar. Nach Ablauf der maximale Dauer wird die Verbindung getrennt, was zu einem Netzwerkfehler zwischen den beteiligten Komponenten führt.

    \item \textbf{Server-Simulator:} Der Server-Simulator simuliert die nachgelagerten Services \gls{srv}~2 und \gls{srv}~3 und stellt damit den Endpunkt der experimentellen Verarbeitungskette dar. Der Server-Simulator empfängt Alarmmeldungen von \gls{lps} bzw.\,w.\ \gls{lptwo} und antwortet anhand von einer stellbaren willkürlichen Fehlerrate mit einer positiven oder negativen Bestätigung. Dadurch bildet der Server-Simulator echte Verarbeitungsfehler in den nachgelagerten Services \gls{srv}~2 und \gls{srv}~3. Außerdem ist die willkürliche Verarbeitungsdauer einer Alarmmeldung einstellbar. Dadurch simuliert der Server beliebige Verarbeitungszeiten der nachgelagerten Services \gls{srv}~2 und \gls{srv}~3. Der Server-Simulator ist außerdem mit einer einstellbaren Idle-Zeit konfigurierbar. Die Idle-Zeit dient dazu, der Server-Simulator zu terminieren und somit die Simulation zu beenden, nachdem eine bestimmte Zeit lang keine Alarmmeldungen mehr empfangen wurden. Abschließend werden erfolgreich verarbeitete Alarmmeldungen in einer Datei gespeichert (\texttt{server-results.json}).

    \item \textbf{Validator:} Der Validator wird nach Abschluss einer Simulation ausgeführt und vergleicht die vom Alarmanlagen-Simulator gesendeter Alarmmeldungen mit den beim Server-Simulator angekommenen Alarmmeldungen. Konkret werden die Vollständigkeit der Zustellung der Alarmmeldungen sowie die Einhaltung der \textit{In-Order-Delivery} und \textit{Exactly-Once-Delivery} validiert. Die Vollständigkeit wird validiert, in dem geprüft wird, dass jede gesendete Alarmmeldung aus (\texttt{device-results.json}) mindestens einmal in (\texttt{server-results.json}) vorkommt. \textit{Exactly-Once-Delivery} wird dadurch validiert, dass jede gesendete Alarmmeldung aus (\texttt{device-results.json}) höchstens einmal in (\texttt{server-results.json}) vorkommt. Um \textit{In-Order-Delivery} zu validieren, werden zunächst mehrfach angekommene Alarmmeldungen entfernt und danach wird geprüft, ob alle Alarmmeldungen jeder einzelnen Alarmanlage anhand von \texttt{triggertime} monoton steigend sind. Jede Verletzung der monotonen Steigfähigkeit einer Alarmanlage wird als \textit{Außer Reihenfolge} klassifiziert. Zur Veranschaulichung der Validierung von \textit{In-Order-Delivery} wird die folgende Sequenz empfangener Alarmmeldungen betrachtet:

          \begin{table}[h]
              \centering
              \begin{tabular}{@{}ccc@{}}
                  \toprule
                  \textbf{Index} & \textbf{device\_id} & \textbf{triggertime} \\
                  \midrule
                  \#1            & 1                   & 0                    \\
                  \#2            & 1                   & 2                    \\
                  \#3            & 1                   & 3                    \\
                  \#4            & 1                   & 1                    \\
                  \#5            & 1                   & 4                    \\
                  \bottomrule
              \end{tabular}
              \caption{Beispielhafte Sequenz empfangener Alarmmeldungen zur Veranschaulichung von In-Order-Delivery Validierung}
              \label{tab:inorder-delivery}
          \end{table}
\end{itemize}


% Da die Experimente stochastisch erfolgen, werden sie solange wiederholt, bis eine Verletzung der Garantie auftritt.

% 3. Vendor lock in

% TODO: for ALPS diagram. We replaced the big diagram shown on page 49 with 2 smaller ones: alps_normal_sequence and alps_retry_sequence



\section{Abgeleitete Anforderungen}
\label{sec:requirements}
% Review-Punkt 8
Aus den in Abschnitt~\ref{sec:mqtt-proxy-schwach} identifizierten Limitationen lassen sich konkrete Anforderungen an \gls{alps} ableiten. Die folgende Liste bündelt diese Anforderungen in einer Form, die direkt in Architekturentscheidungen überführt und in der Evaluation überprüft werden kann.

\begin{itemize}
    \item \textbf{A1 (Asynchrone Entkopplung beibehalten):} Vorgelagerte Komponenten dürfen nicht auf die vollständige Verarbeitung in allen nachgelagerten Services warten müssen. Eingehende Meldungen sollen daher weiterhin brokerbasiert entgegengenommen und zwischengespeichert werden.

    \item \textbf{A2 (In-Order-Delivery pro Alarmanlage sicherstellen):} Meldungen derselben Alarmanlage müssen in der Reihenfolge verarbeitet werden, in der sie erzeugt wurden. Diese Eigenschaft muss auch bei temporären Fehlern und Wiederholungsverarbeitung erhalten bleiben.

    \item \textbf{A3 (Fachliche Duplikate robust behandeln):} Das Design muss wiederholte Übertragungen derselben fachlichen Meldung erkennen oder ihre Auswirkungen zuverlässig begrenzen. Voraussetzung ist eine eindeutige fachliche Identifikation pro Meldung.

    \item \textbf{A4 (Fehlerfälle ohne Nachrichtenverlust behandeln):} Netzwerkfehler, negative Antworten nachgelagerter Services und temporäre Ausfälle dürfen nicht zu Nachrichtenverlust führen. Fehlgeschlagene Meldungen müssen erneut verarbeitet werden können, ohne die Reihenfolgegarantie pro Alarm\-an\-la\-ge zu verletzen.

    \item \textbf{A5 (Kontrollierte Parallelisierung ermöglichen):} Das System soll mehrere Alarmanlagen parallel verarbeiten können. Gleichzeitig muss die Verarbeitung innerhalb einer einzelnen Alarmanlage geordnet bleiben.

    \item \textbf{A6 (Betriebsaufwand bei Skalierung begrenzen):} Zusätzliche Last soll nicht primär durch beliebig viele identische Service-Instanzen absorbiert werden. Das Design soll Parallelität gezielt über Broker-, Queue- oder Partitionierungsmechanismen abbilden, um die operative Komplexität zu reduzieren.

    \item \textbf{A7 (Vendor-Lock-in reduzieren):} Das Design soll möglichst auf standardnahen und portablen Mechanismen aufbauen. Falls plattformspezifische Funktionen erforderlich sind, müssen sie klar isoliert und mit vertretbarem Aufwand austauschbar sein.

    \item \textbf{A8 (Transparenz zum dem Legacy-System bewahren):} \gls{alps} darf die fachliche Schnittstelle des Legacy-Systems nicht verändern. Vor- und nachgelagerte Komponenten sollen weiterhin über ihre bestehenden Kommunikationsschnittstellen integrierbar bleiben.
\end{itemize}

Diese Anforderungen bilden die Grundlage für das im folgenden Kapitel vorgestellte \gls{amqp}-basierte \gls{alps}.

\chapter{\acrfull{alps}}
\label{ch:amqp-proxy}

Das in diesem Kapitel beschriebene \gls{amqp}-basierte \gls{alps} wird als Weiterentwicklung des zuvor betrachteten \gls{mqtt}-basierten \gls{lps} verstanden. Die grundlegende Idee der asynchronen Entkopplung bleibt erhalten: Vorgelagerte Komponenten sollen weiterhin nicht auf die vollständige Verarbeitung durch alle nachgelagerten Services warten müssen. Gleichzeitig adressiert das neue Design die im vorherigen Kapitel identifizierten Limitationen des bestehenden Ansatzes. Dazu gehören eine besser kontrollierbare Skalierung, eine zuverlässigere Behandlung von Fehlerfällen, die Sicherstellung der Reihenfolge pro Alarmanlage sowie eine robuste Behandlung von Duplikaten.

Die Gesamtidee des Systems besteht darin, die bisherige Verarbeitungskette in vier Teile aufzuteilen: einen Producer-Teil (\gls{alpone}), einen Broker, einen Consumer-Teil (\gls{alptwo}) und einen zusätzlichen Service zur Duplikaterkennung (\gls{alpthree}). \gls{alpone} übernimmt die Rolle von \gls{lpone}. % Review-Punkt 53: einheitlich HTTP-Requests
% Review-Punkt 54/61: Event = Request-Nutzlast + Metadaten (AMQP-Terminologie)
Er nimmt \gls{http}-Requests von vorgelagerten Komponenten entgegen, wandelt sie in interne Events um\footnote{Unter einem internen Event wird die Originalnachricht des vorgelagerten Services verstanden, die um Metadaten wie ID und Zeitstempel ergänzt wurde. Anhand dieser Daten lässt sich jede Alarmmeldung eindeutig identifizieren, was für den nachgelagerten Deduplizierungsschritt in \gls{alpthree} erforderlich ist.} und publiziert sie an den \gls{amqp}-Broker. \gls{alptwo} übernimmt die Rolle von \gls{lptwo}. Er konsumiert die Events aus dem Broker und leitet die ursprüngliche Nachricht nicht direkt an \gls{srv}~3 weiter, sondern kommuniziert mit \gls{alpthree} zur Duplikaterkennung. \gls{alpthree} implementiert die Schnittstelle von \gls{srv}~3 und wird vor \gls{srv}~3 geschaltet\footnote{\gls{alpthree} und \gls{srv}~3: Bei $N$ Services würde \gls{alpthree} unmittelbar vor dem $N$-ten (letzten) Service in der Kette bereitgestellt werden. Die Zahl 3 resultiert daraus, dass \gls{srv}~3 im originalen Paper so genannt wurde und durchgängig den letzten Service in der Kette repräsentiert. Daher wurde analog der Name \gls{alpthree} gewählt.}. Er erkennt beziehungsweise vermeidet Duplikate und leitet die Nachricht anschließend an \gls{srv}~3 weiter.

% Review-Punkt 55: triggertime als eindeutige Auslösezeit präzisieren
Im Folgenden wird angenommen, dass Alarmmeldungen eine \gls{json}-Struktur mit mindestens den Feldern \texttt{device\_id} und \texttt{triggertime} besitzen. \newline Die \texttt{device\_id} ist eindeutig pro Alarmanlage und wird verwendet, um die Anlage zu identifizieren. Das Feld \texttt{triggertime} ist die eindeutige Auslösezeit einer Alarmmeldung einer Alarmanlage und dient dazu, eine einzelne fachliche Meldung zu identifizieren. Wiederholte Meldungen besitzen dabei dieselbe \texttt{triggertime} wie die ursprüngliche Meldung.

\section{Designziele}

Die Designziele von \gls{alps} ergeben sich direkt aus den Limitationen von \gls{lps}. Das neue Design soll die asynchrone Entkopplung beibehalten, aber die Verarbeitung stärker an den fachlichen Anforderungen des Legacy-Systems ausrichten. Zentral ist dabei, dass Nachrichten einer einzelnen Alarmanlage beziehungsweise einer einzelnen Alarmanlage in ihrer ursprünglichen Reihenfolge verarbeitet werden müssen, während Nachrichten unterschiedlicher Alarmanlagen parallel verarbeitet werden dürfen. Gleichzeitig soll die bestehende Legacy-Schnittstelle transparent weiter nutzbar bleiben.

Aus diesen Vorgaben ergeben sich die folgenden Ziele:
\begin{itemize}
    \item Die asynchrone Entkopplung der ursprünglichen Request-Response-Kette soll beibehalten werden.
    \item Die Verarbeitung soll \textit{In-Order-Delivery} pro Alarmanlage sicherstellen.
    \item Die Skalierbarkeit soll mit geringerem operativem Aufwand ermöglicht werden.
    \item Fehlerfälle wie Serverfehler und Netzwerkfehler sollen sauber behandelt werden.
    \item Duplikate sollen erkannt oder ihre Auswirkungen begrenzt werden.
    \item Vendor-Lock-in soll reduziert werden, indem möglichst standardnahe \gls{amqp}-Mechanismen verwendet werden, die bei mehreren Brokern in ähnlicher Form verfügbar sind.
\end{itemize}

\section{Architektur und Komponenten}

% Review-Punkt 57: Anschluss analog zum originalen LPS
\gls{alps} besteht aus vier zentralen Komponenten: \gls{alpone}, dem \gls{amqp}-Broker, \gls{alptwo} und \gls{alpthree} zur Duplikaterkennung. Diese Aufteilung übernimmt die Grundstruktur des zuvor betrachteten \gls{lps}, ersetzt jedoch die \gls{mqtt}-basierte Kommunikation durch eine \gls{amqp}-basierte Broker-Struktur und ergänzt die Verarbeitungskette um eine Komponente zur Behandlung fachlicher Duplikate. \gls{alpone} entspricht funktional \gls{lpone}, während \gls{alptwo} die Rolle von \gls{lptwo} übernimmt, analog zum originalen \gls{lps}. Daraus folgt, dass die asynchrone Entkopplung zwischen vorgelagertem und nachgelagertem Servicepfad erhalten bleibt und die bestehenden Legacy-Schnittstellen transparent weiter nutzbar sind. Zwischen \gls{alpone} und \gls{alptwo} liegt der Broker, der eingehende Events entgegennimmt, zwischenspeichert und anhand des verwendeten Rout\-ing-Me\-cha\-nis\-mus auf Queues verteilt. Hinter \gls{alptwo} wird der zusätzliche \gls{alpthree} vor \gls{srv}~3 geschaltet.

\gls{alpone} bildet die Eingangskomponente des Systems. Er empfängt eingehende Requests von vorgelagerten Komponenten über die bestehende \gls{http}-Schnittstelle. Anschließend wandelt er die empfangene Request in ein internes Event um und publiziert dieses Event an den \gls{amqp}-Broker. Nach erfolgreicher Übergabe an den Broker kann \gls{alpone} eine Bestätigung an die vorgelagerte Komponente zurückgeben. Dadurch wird die ursprüngliche synchrone Request-Response-Kette an dieser Stelle entkoppelt, analog zum originalen \gls{lps}: Die vorgelagerte Komponente muss nicht warten, bis die nachgelagerte Verarbeitung vollständig abgeschlossen ist.

Für das Publishing verbindet sich \gls{alpone} mit dem Broker und deklariert den Exchange \texttt{events.hash} vom Typ \texttt{x-consistent-hash}. Events werden mit der \texttt{device\_id} als Routing Key veröffentlicht. Der Routing Key repräsentiert dabei die Alarmanlage, zu der die Meldung gehört. Der Broker entscheidet auf dieser Grundlage, in welche Queue die Nachricht geschrieben wird. Nachrichten mit derselben \texttt{device\_id} sollen dadurch derselben Queue zugeordnet werden, während Nachrichten unterschiedlicher Alarmanlagen auf verschiedene Queues verteilt werden können.

Beim Start legt \gls{alpone} mehrere Queues an und bindet diese Queues an den Consistent-Hash-Exchange. Jede Queue wird mit dem Gewicht \texttt{1} an den Exchange gebunden. Zusätzlich wird für die Queues \texttt{x-single-active-consumer} gesetzt. Dadurch wird festgelegt, dass pro Queue jeweils nur ein aktiver Consumer Nachrichten verarbeitet, auch wenn mehrere Consumer für dieselbe Queue registriert sind. Die Anzahl der Queues bildet damit einen zentralen Parameter für die mögliche Parallelisierung zwischen unterschiedlichen Alarmanlagen.

\gls{alptwo} bildet die Ausgangskomponente der brokerbasierten Verarbeitung. Er konsumiert Events aus dem Broker und sendet die ursprüngliche Nachricht an \gls{alpthree} zur Duplikaterkennung weiter. % Review-Punkt 60: Task als technische Arbeitseinheit/Consumer-Loop definieren
Die Implementierung von \gls{alptwo} startet für jede Queue eine eigene \textit{Task}\footnote{Task bezeichnet hier eine dedizierte Verarbeitungsroutine beziehungsweise Consumer-Schleife von \gls{alptwo}, die genau einer Queue zugeordnet ist.}. Jede Task konsumiert aus genau einer Queue und arbeitet mit \texttt{prefetch\_count = 1}. Dadurch wird pro Task immer nur eine Nachricht gleichzeitig an den \gls{alptwo} ausgeliefert. Diese Beschränkung reduziert die Parallelität innerhalb einer Queue bewusst, schützt aber die geordnete Verarbeitung der Nachrichten, die dieser Queue zugeordnet sind.

% Review-Punkt 62: explizit AMQP-ACK und AMQP-NACK
Das Bestätigungsverhalten von \gls{alptwo} ist an die erfolgreiche Verarbeitung durch den nachgelagerten Service gekoppelt. Eine Nachricht wird erst dann mit einem \gls{amqp}-\gls{ack} bestätigt, wenn der \gls{http}-Request an diesen Service erfolgreich war. Tritt ein Fehler auf, beispielsweise durch eine negative Antwort oder einen Netzwerkfehler, bestätigt \gls{alptwo} die Nachricht nicht erfolgreich. Stattdessen wird ein \gls{amqp}-\gls{nack} mit \texttt{requeue=true} gesendet, damit der Broker die Nachricht erneut zur Verarbeitung bereitstellt. Anschließend wird die \gls{amqp}-Verbindung geschlossen und neu aufgebaut. Wenn die Verbindung geschlossen ist, kann der Broker die betroffene Queue einem anderen Task in einer anderen Instanz von \gls{alptwo} zuweisen. Dadurch kann eine erneut bereitgestellte Nachricht von einer anderen Instanz von \gls{alptwo} verarbeitet werden, ohne dass sie verloren geht. Gleichzeitig unterstützt dieses Verhalten eine gleichmäßigere Verteilung von Queues beziehungsweise Meldungen auf Tasks in verschiedenen \gls{alptwo}-Instanzen.

\gls{alpthree} ergänzt die Architektur um eine dedizierte Stelle zur Behandlung von Duplikaten. Er implementiert nach außen die Schnittstelle von \gls{srv}~3 und kann dadurch vor \gls{srv}~3 geschaltet werden, ohne dass \gls{alptwo} direkt mit \gls{srv}~3 kommunizieren muss. Aus Sicht von \gls{alptwo} ist \gls{alpthree} der nachgelagerte Kommunikationspartner. Aus Sicht von \gls{srv}~3 übernimmt \gls{alpthree} die Rolle einer vorgelagerten Schutzschicht: Er prüft eingehende Meldungen anhand von \texttt{device\_id} und \texttt{triggertime} auf wiederholte fachliche Nachrichten und leitet nur die eindeutigen Nachrichten weiter, die an \gls{srv}~3 übergeben werden sollen.

% Review-Punkt 22: aktive Abbildungsbezüge
Abbildung~\ref{fig:amqp_proxy_architecture} zeigt die Gesamtarchitektur des entwickelten \gls{alps} mit den Komponenten \gls{alpone}, Broker, \gls{alptwo} und der Deduplizierungsschicht \gls{alpthree} vor \gls{srv}~3.

\begin{figure}[hbtp]
    \centering
    \includegraphics[width=0.95\textwidth]{./img/a_lps.pdf}
    \caption{Gesamtarchitektur des entwickelten \acrshort{amqp}-basierten \acrshort{alps}.}
    \label{fig:amqp_proxy_architecture}
\end{figure}

\section{\acrshort{amqp}-Mechanismen im Design}

Als \gls{amqp}-Broker wird für die konkrete Umsetzung von \gls{alps} \textit{RabbitMQ}~\citep{rabbitmq} eingesetzt. Wie bereits in Abschnitt~\ref{subsec:grundlagen_amqp} dargelegt, sind die verwendeten Konfigurationen und Architekturkonzepte jedoch nicht exklusiv an RabbitMQ gebunden. Die eingesetzten Mechanismen lassen sich funktional auch auf andere gängige Message Broker, wie beispielsweise \textit{Apache ActiveMQ Artemis} oder \textit{LavinMQ}, übertragen. \gls{alps} verwendet mehrere standardnahe \gls{amqp}- sowie brokerspezifische Mechanismen, um Routing, Zustellkontrolle und Fehlerbehandlung miteinander zu verbinden.

Die zentralen Bausteine sind \texttt{x-consistent-hash}, \texttt{prefetch\_count}, \newline \texttt{x-single-active-consumer}, mehrere Queues, Consumer Acknowledgments  sowie  Negative Acknowledgments mit Requeue.

Für das Routing wird ein Exchange vom Typ \texttt{x-consistent-hash} eingesetzt. Dieser verteilt Nachrichten anhand eines stabilen Schlüssels deterministisch auf Queues. In Verbindung mit mehreren Queues wird damit eine kontrollierte Parallelisierung zwischen unterschiedlichen Alarmmeldungen ermöglicht. Es wird garantiert, dass Alarmmeldungen einer Alarmanlage immer derselben Queue zugewiesen werden.

% Review-Punkt 62: explizit AMQP-ACK und AMQP-NACK
Für die Zustellkontrolle wird das Acknowledgment-Modell des Brokers verwendet. Nachrichten gelten erst dann als abgeschlossen, wenn ein Consumer sie nach erfolgreicher nachgelagerter Verarbeitung mit einem \gls{amqp}-\gls{ack} bestätigt. Im Fehlerfall kann ein Consumer mit einem \gls{amqp}-\gls{nack} und \texttt{requeue=true} eine erneute Bereitstellung der Nachricht auslösen.

Zusätzlich begrenzen \texttt{prefetch\_count = 1} und \texttt{x-single-active-consumer} die gleichzeitige Verarbeitung pro Queue. Dadurch wird pro Queue jeweils nur eine Nachricht aktiv verarbeitet, während bei Ausfall eines Consumers eine Übernahme durch einen anderen Consumer möglich bleibt.

Wie diese Mechanismen im konkreten Architekturentwurf zur Sicherstellung der \textit{In-Order-Delivery} eingesetzt werden, beschreibt der folgende Abschnitt.

% Review-Punkt 22: aktive Abbildungsbezüge
Abbildung~\ref{fig:amqp_mechanisms_} zeigt die zentralen \gls{amqp}-Mechanismen des Designs, insbesondere Consistent-Hash-Routing, Queue-Bindings, \texttt{x-single-active\-consumer}, \texttt{prefetch\_count} und das Ack/Nack-Verhalten.

\begin{figure}[hbtp]
    \centering
    \includegraphics[width=0.95\textwidth]{./img/a_lps_intern.pdf}
    \caption{Interne \acrshort{amqp}-Mechanismen von \acrshort{alps} (Exchange-, Queue- und Consumer-Interaktion).}
    \label{fig:amqp_mechanisms_}
\end{figure}

\section{Sicherstellung von \textit{In-Order-Delivery}}
Die angestrebte \textit{In-Order-Delivery} wird in \gls{alps} als Reihenfolgegarantie pro Alarmanlage umgesetzt. Dazu werden alle Nachrichten einer Alarmanlage mit derselben \texttt{device\_id} als Routing Key publiziert. Der Consistent-Hash-Exchange ordnet identische Schlüssel deterministisch derselben Queue zu, sodass alle Nachrichten einer Alarmanlage entlang derselben Queue-Sequenz verarbeitet werden.

Innerhalb dieser Queue wird die Reihenfolge über die Kombination von \texttt{x-single-active-consumer} und \texttt{prefetch\_count = 1} abgesichert. Pro Queue ist jeweils nur ein Consumer aktiv, und pro Consumer wird jeweils nur eine unbestätigte Nachricht gleichzeitig verarbeitet. Dadurch kann keine spätere Nachricht derselben Alarmanlage parallel an einer früheren Nachricht vorbeilaufen.

% Review-Punkt 62: explizit AMQP-ACK und AMQP-NACK
Im Fehlerfall wird eine Nachricht nicht mit einem \gls{amqp}-\gls{ack} abgeschlossen, sondern per \gls{amqp}-\gls{nack} mit \texttt{requeue=true} erneut zur Verarbeitung bereitgestellt. Damit bleibt eine fehlgeschlagene frühere Nachricht der limitierende Schritt innerhalb der betreffenden Gerätefolge. Wenn eine Nachricht wiederholt fehlschlägt, entsteht für diese Queue ein Engpass, da nachfolgende Nachrichten derselben Alarmanlage erst nach erfolgreichem Abschluss der vorherigen Nachricht fachlich korrekt weiterverarbeitet werden können.

Das Design zielt nicht auf eine globale Reihenfolge über alle Alarmanlagen hinweg ab. Garantiert beziehungsweise angestrebt wird die Reihenfolge pro Alarmanlage, während Nachrichten unterschiedlicher Alarmanlagen parallel über verschiedene Queues verarbeitet werden dürfen.
% Review-Punkt 22: aktive Abbildungsbezüge
Abbildung~\ref{fig:amqp_in_order_delivery_placeholder} zeigt, wie Nachrichten pro \texttt{device\_id} deterministisch einer Queue zugeordnet und dort sequentiell verarbeitet werden, sodass die Reihenfolge je Alarmanlage erhalten bleibt.

\begin{figure}[hbtp]
    \centering
    \includegraphics[width=0.95\textwidth]{./diagram/a_lps_in_order_delivery.png}
    \caption{Mechanismus zur Sicherstellung der \textit{In-Order-Delivery} pro Alarmanlage in \acrshort{alps}.}
    \label{fig:amqp_in_order_delivery_placeholder}
\end{figure}

\section{Sicherstellung von \textit{Exactly-Once-Delivery}}

Die zuvor beschriebenen \gls{amqp}-Mechanismen stellen eine zuverlässige Zustellung im Sinne erneuter Verarbeitung nicht bestätigter Nachrichten bereit. Für sich allein garantieren sie jedoch keine fachliche \textit{Exactly-Once-Delivery} über die gesamte Verarbeitungskette hinweg. Der Grund liegt darin, dass der Broker nur den Zustand der Nachricht innerhalb des \gls{amqp}-Systems kennt. Ob eine Nachricht bereits an den nachgelagerten Service übergeben und dort verarbeitet wurde, kann der Broker nicht zuverlässig beurteilen, wenn während der Verarbeitung ein Fehler in der Verbindung zum Consumer auftritt.

Ein typischer Fehlerfall verdeutlicht dieses Problem. Ein Consumer entnimmt eine Nachricht aus dem Broker und sendet sie an den nachgelagerten Service. Während der Consumer auf die Antwort wartet, bricht die Verbindung zwischen Consumer und Broker ab. % Review-Punkt 62: explizit AMQP-ACK und AMQP-NACK
% Review-Punkt 64: klar zwei verschiedene Consumer-Instanzen bei Re-Delivery
Da der Broker kein \gls{amqp}-\gls{ack} erhält, behandelt er die Nachricht als nicht erfolgreich verarbeitet und stellt sie erneut zu. Eine andere Consumer-Instanz (nicht dieselbe Prozessinstanz) kann dieselbe Nachricht anschließend erneut an den nachgelagerten Service senden. Aus Sicht des Brokers ist dieses Verhalten korrekt, weil es Nachrichtenverlust verhindert. Aus Sicht des Legacy-Systems entsteht jedoch ein fachliches Duplikat, da dieselbe Alarmmeldung mehrfach bei \gls{srv}~3 ankommen kann.

\begin{figure}[hbtp]
    \centering
    \includegraphics[width=0.95\textwidth]{./diagram/a_lps_exactly_once_delivery_dup.png}
    % Review-Punkt 62
    \caption{Fehlersequenz zur Entstehung fachlicher Duplikate trotz AMQP-ACK-after-processing.}
    \label{fig:amqp_duplicate_sequence}
\end{figure}
% Review-Punkt 22: aktive Abbildungsbezüge
% Review-Punkt 62: explizit AMQP-ACK und AMQP-NACK
Abbildung~\ref{fig:amqp_duplicate_sequence} zeigt den kritischen Fehlerpfad: Nach erfolgreicher Weitergabe an den nachgelagerten Service geht das \gls{amqp}-\gls{ack} aufgrund eines Verbindungsabbruchs verloren, wodurch der Broker dieselbe Nachricht erneut zustellt.

Zur Behandlung dieser Problematik wird die \gls{amqp}-basierte Zustelllogik um \gls{alpthree} ergänzt. \gls{alpthree} ist nicht Teil des Brokers und ersetzt auch nicht die Consumer-Acknowledgments. Stattdessen bildet er eine fachliche Deduplizierungsschicht unmittelbar vor \gls{srv}~3. \gls{alptwo}-Instanzen kommunizieren daher nicht direkt mit \gls{srv}~3, sondern mit \gls{alpthree}. \gls{alpthree} implementiert die Schnittstelle von \gls{srv}~3, nimmt die vom \gls{alptwo} weitergeleitete Alarmmeldung entgegen, prüft ihre fachliche Identität und entscheidet anschließend, ob die Meldung an \gls{srv}~3 weitergegeben wird.

Die Deduplizierung basiert auf der Annahme, dass wiederholte Übertragun\-gen derselben fachlichen Alarmmeldung dieselbe Anlagenkennung und dieselbe \newline \texttt{triggertime} besitzen. Diese beiden Informationen bilden zusammen die fachliche Identität einer Meldung. Geht eine Meldung bei \gls{alpthree} ein, wird geprüft, ob diese fachliche Identität bereits als verarbeitet bekannt ist. Ist dies der Fall, wird die Meldung als Duplikat behandelt und nicht erneut an \gls{srv}~3 weitergeleitet. \gls{alpthree} kann \gls{alptwo} dennoch eine erfolgreiche Antwort geben, da die Meldung aus fachlicher Sicht bereits berücksichtigt wurde. Ist die fachliche Identität noch nicht bekannt, wird die Meldung an \gls{srv}~3 weitergeleitet und anschließend als verarbeitet registriert.

% Review-Punkt 22: aktive Abbildungsbezüge
Abbildung~\ref{fig:a_lps_srv4} zeigt die Einbettung von \gls{alpthree} als dedizierte Deduplizierungsschicht vor \gls{srv}~3. Dadurch werden fachliche Duplikate vor der Übergabe an den Legacy-Service erkannt und gefiltert.

\begin{figure}[hbtp]
    \centering
    \includegraphics[width=0.95\textwidth]{./img/a_lps_3.pdf}
    \caption{\acrshort{alpthree} als vorgelagerte Deduplizierungsschicht vor \acrshort{srv}~3.}
    \label{fig:a_lps_srv4}
\end{figure}

\subsection{Mnesia}
Für die Implementierung von \gls{alpthree} wird das verteilte Datenbanksystem Mnesia im Cluster-Modus eingesetzt. Innerhalb der Datenbank wird eine verteilte Tabelle mit dem Speichertyp \texttt{\textbf{disc\_copies}} angelegt (vgl. Abschnitt~\ref{subsec:grundlagen_mnesia}). Dies gewährleistet, dass die Daten über mehrere Knoten hinweg repliziert und persistent auf dem Festspeicher abgelegt werden, sodass der Systemzustand auch bei Hardwareausfällen oder Neustarts einzelner Knoten verlustfrei erhalten bleibt.

Zur zuverlässigen Identifikation von Duplikaten wird in dieser Tabelle die Konkatenation aus \texttt{device\_id} und \texttt{triggertime} gespeichert. Diese Konkatenation identifiziert jede fachliche Alarmmeldung eindeutig. Der prozedurale Ablauf stellt dabei sicher, dass ein neuer Eintrag erst dann in die Datenbank geschrieben wird, nachdem die entsprechende Meldung vom nachgelagerten \gls{srv}~3 erfolgreich verarbeitet wurde.

Um der Problematik divergierender Datenbestände bei temporären Netzwerkpartitionierungen (Split-Brain-Szenarien; vgl. Abschnitt~\ref{subsec:grundlagen_mnesia}) zu begegnen, % Review-Punkt 65: „basierend auf“ statt „anhand“/„mittels“-Doppelung
ist das Cluster basierend auf der Bibliothek \textit{unsplit}~\citep{unsplit2010} so konfiguriert, dass nach der Wiederherstellung der Netzwerkverbindung die Einträge aller isolierten Knoten automatisch zusammengeführt werden. Diese Strategie der Konfliktauflösung ist in diesem spezifischen Anwendungsfall sicher und unproblematisch anwendbar, da bestehende Datensätze in der Datenbank niemals nachträglich modifiziert oder gelöscht werden. Die Datenbasis wird stattdessen ausschließlich durch das Hinzufügen neuer Einträge nach einer erfolgreichen Verarbeitung durch \gls{srv}~3 erweitert (\textit{Append-Only-Prinzip}\footnote{Append-Only bedeutet hier, dass bestehende Einträge nicht geändert oder gelöscht, sondern nur neue Einträge ergänzt werden.}). Dadurch wird garantiert, dass alle Knoten nach dem Zusammenführen der Daten wieder über einen korrekten, konsistenten Gesamtzustand verfügen und den regulären Betrieb nahtlos fortsetzen können.

% Review-Punkt 66: Seitenumbruch vor Deployment, um hängenden Einzelsatz zu vermeiden
\pagebreak
\subsection{Deployment}
Jede Instanz von \gls{alpthree} wird nach dem \textit{Sidecar-Prinzip}\footnote{Das Sidecar-Prinzip beschreibt die Bereitstellung einer unterstützenden Komponente in unmittelbarer Nähe zu einer Hauptkomponente, sodass beide eng gekoppelt zusammenarbeiten.}~\citep{burns2018designing} unmittelbar vor \gls{srv}~3 bereitgestellt. Diese direkte Nähe zu \gls{srv}~3 ist wichtig für eine zuverlässige \textit{Exactly-Once-Delivery}. Je länger der Netzwerkpfad zwischen der Deduplizierung und \gls{srv}~3 ist, desto größer ist das Risiko für Kommunikationsfehler. Ein Netzwerkfehler oder Timeout könnte genau dann auftreten, nachdem \gls{srv}~3 die Meldung bereits erfolgreich verarbeitet hat. \gls{alpthree} erhält in diesem Fall keine Bestätigung von \gls{srv}~3. Dadurch wird die Alarmmeldung in der Datenbank nicht als verarbeitet markiert. Bei einem erneuten Sendeversuch wird die Meldung dann ein zweites Mal an \gls{srv}~3 übergeben, wodurch ein fachliches Duplikat entsteht. Die direkte Positionierung von \gls{alpthree} als vorgelagerte Sidecar-Komponente minimiert genau diesen kritischen Kommunikationspfad.
Zusammenfassend verbindet \gls{alps} damit asynchrone Entkopplung, \textit{In-Order-Delivery} pro Alarmanlage und fachliche Deduplizierung in einer Architektur, die ohne Änderungen am Legacy-Quellcode integrierbar bleibt. Auf dieser Grundlage wird im folgenden Kapitel experimentell überprüft, inwieweit die abgeleiteten Anforderungen unter Last- und Fehlerszenarien erfüllt werden.

\chapter{Experimentelle Evaluation}
\label{ch:eval}

Nachdem im vorherigen Kapitel Design und Implementierung von \gls{alps} vorgestellt wurden, überprüft dieses Kapitel experimentell, ob die abgeleiteten Anforderungen (vgl. Abschnitt~\ref{sec:requirements}) tatsächlich erfüllt werden. Zwei Fragen stehen dabei im Vordergrund: Erstens, ob \gls{alps} die zentralen Delivery-Garantien \textit{Exactly-Once-Delivery} und \textit{In-Order-Delivery} auch unter Fehlerbedingungen einhält. Zweitens, wie sich das System bei steigender Last verhält und über welchen Mechanismus zusätzliche Kapazität am effizientesten bereitgestellt werden kann. Um beide Fragen sauber voneinander zu trennen, ist die Evaluation in zwei Phasen mit jeweils eigenem Schwerpunkt gegliedert:

\begin{description}
    \item[Phase 1 -- Korrektheit unter Fehlerbedingungen:]
          Hier wird geprüft, ob \gls{alps} unter denselben Fehlerbedingungen (Netzwerkfehler und Serverfehlerantworten), an denen das ursprüngliche \gls{lps} in Abschnitt~\ref{sub-sub-sec:phase-2} gescheitert ist, die Delivery-Garantien einhält.

    \item[Phase 2 -- Skalierbarkeit im Normalbetrieb:]
          Hier wird im fehlerfreien Referenzbetrieb untersucht, ob sich zusätzliche Last eher durch mehr Queues oder durch mehr \gls{alptwo}-Instanzen abfedern lässt, und was das für den operativen Aufwand bedeutet.
\end{description}

Damit die Ergebnisse beider Architekturvarianten direkt vergleichbar bleiben, verwendet die Durchführung dieselben Hilfsservices wie die Evaluation des \gls{lps} in Abschnitt~\ref{sec:mqtt-proxy-experimente}: Alarmanlagen-Simulator, \gls{tcp}-Proxy, Server-Simulator und Validator. Alle Experimente laufen mit Docker Compose~\citep{docker_compose}; die verwendeten Konfigurationen sowie Hinweise zur Reproduktion sind in Abschnitt~\ref{sec:appendix-reproducibility} dokumentiert. Als Testsystem diente ein Rechner mit einer \textit{Intel(R) Core(TM) i5-9600K} CPU~\citep{intel_i5_9600k}, 16\,GB RAM und \textit{Ubuntu} 24.04.4~\citep{ubuntu_24_04}.

\section{Phase 1: Korrektheit unter Fehlerbedingungen}
\label{sec:eval-korrektheit}

Ziel dieser Phase ist der Korrektheitsnachweis von \gls{alps} unter Fehlereinfluss. Dazu wird analog zu Phase~2 der \gls{lps}-Evaluation (Abschnitt~\ref{sub-sub-sec:phase-2}) eine Kombination aus Netzwerkfehlern (\gls{tcp}-Drops zwischen Broker und \gls{alptwo}) und Serverfehlerantworten (fehlerhafte Antworten von \gls{srv}~3) injiziert. Weil dasselbe Fehlerprofil verwendet wird, das beim \gls{lps} zu 77{,}8\,\% Duplikaten und 13{,}5\,\% Reihenfolgeverletzungen geführt hat, lässt sich unmittelbar beurteilen, ob \gls{alps} diese Schwachstellen tatsächlich behebt.

Tabelle~\ref{tab:amqp_exp_phase_1_config} zeigt die verwendete Konfiguration. Sie orientiert sich an der Größenordnung der \gls{lps}-Baseline, ergänzt um zehn Queues sowie drei \gls{alpthree}-Instanzen zur Duplikaterkennung.

\begin{table}[H]
    \centering
    \caption{Konfiguration des Korrektheitsexperiments in Phase~1.}
    \label{tab:amqp_exp_phase_1_config}
    \begin{tabular}{ll}
        \hline
        Parameter                   & Wert      \\
        \hline
        Simulierte Alarmanlagen     & 100       \\
        Meldungen je Alarmanlage    & 100       \\
        \gls{alpone}-Instanzen      & 3         \\
        \gls{alptwo}-Instanzen      & 3         \\
        Queues                      & 10        \\
        \gls{tcp}-Proxy-Fehlerrate  & 0{,}1\,\% \\
        Serverfehlerrate            & 0{,}1\,\% \\
        Maximale Server-Antwortzeit & 100\,ms   \\
        \gls{alpthree}-Instanzen    & 3         \\
        \hline
    \end{tabular}
\end{table}

Wie schon bei der \gls{lps}-Evaluation erfolgt die Fehlerinjektion stochastisch: Der Server-Simulator entscheidet pro Anfrage probabilistisch über eine Fehlerantwort, der \gls{tcp}-Proxy probabilistisch über einen Verbindungsabbruch und trennt akzeptierte Verbindungen zusätzlich nach einer zufälligen Dauer. Ausgewertet wird der tatsächlich dokumentierte Lauf dieser Konfiguration.

Abbildung~\ref{fig:amqp_exp_phase_1_duplicate} zeigt die Duplikatprüfung des Validators und damit die \textit{Exactly-Once-Delivery}-Eigenschaft unter den beschriebenen Fehlerbedingungen.

\begin{figure}[H]
    \centering
    \includegraphics[width=1\textwidth]{./img/exp/amqp/1/failure-duplicate.pdf}
    \caption{Validator-Ausgabe zur Duplikatprüfung in der \acrshort{alps}-Evaluation unter Fehlerbedingungen.}
    \label{fig:amqp_exp_phase_1_duplicate}
\end{figure}

Ergänzend zeigt Abbildung~\ref{fig:amqp_exp_phase_1_order} die Reihenfolgeprüfung pro Alarmanlage und damit die \textit{In-Order-Delivery}-Eigenschaft unter denselben Fehlerbedingungen.

\begin{figure}[H]
    \centering
    \includegraphics[width=1\textwidth]{./img/exp/amqp/1/failure-order.pdf}
    \caption{Validator-Ausgabe zur Reihenfolgeprüfung in der \acrshort{alps}-Evaluation unter Fehlerbedingungen.}
    \label{fig:amqp_exp_phase_1_order}
\end{figure}

Beide Auswertungen weisen für die evaluierte Konfiguration keine Verletzung der betrachteten Delivery-Garantien aus: Es treten weder Duplikate noch Reihenfolgeverletzungen auf, und alle erzeugten Meldungen kommen vollständig an. Damit hält \gls{alps} unter genau den Fehlerbedingungen stand, unter denen das \gls{lps} zuvor gescheitert ist, dass die Kombination aus deterministischem Queue-Routing und der nachgelagerten Deduplizierungsschicht \gls{alpthree} wirksam ist. Wie robust dieses Ergebnis über unterschiedliche Lastprofile hinweg ist, wird im folgenden Abschnitt zur Skalierbarkeit weiter eingeordnet.

\section{Phase 2: Skalierbarkeit im Normalbetrieb}
\label{sec:eval-skalierbarkeit}

Nachdem Phase~1 die Korrektheit unter Fehlereinfluss belegt hat, geht es in Phase~2 um die Frage, wie \gls{alps} unter steigender Last im fehlerfreien Referenzbetrieb skaliert. Konkret wird untersucht, ob zusätzliche Kapazität eher durch eine höhere Queue-Anzahl oder durch zusätzliche \gls{alptwo}-Instanzen erreicht wird und welche der beiden Parameter den größeren Effekt auf die Ende-zu-Ende-Latenz hat.

Dazu wird die Versuchsreihe in zwei Blöcke aufgeteilt, die jeweils nur einen Parameter gegenüber einer gemeinsamen Baseline (Phase~2.1) verändern: In Phase~2.2 und Phase~2.3 wird ausschließlich die Queue-Anzahl erhöht, während \gls{alpone}, \gls{alptwo} und das Lastprofil konstant bleiben. In Phase~2.4 und Phase~2.5 wird umgekehrt zur ursprünglichen Queue-Anzahl zurückgekehrt und stattdessen die Anzahl der \gls{alptwo}-Instanzen erhöht. Diese Trennung erlaubt es, den isolierten Effekt beider Parameter direkt gegenüberzustellen. Tabelle~\ref{tab:amqp_exp_scaling_config} fasst die fünf Konfigurationen zusammen; die jeweils veränderten Werte gegenüber der Baseline sind hervorgehoben.

\begin{table}[H]
    \centering
    \caption{Konfigurationen der Skalierungsexperimente für \acrshort{alps}.}
    \label{tab:amqp_exp_scaling_config}
    \begin{tabular}{llllll}
        \hline
        Experiment & \gls{alpone} & Queues      & \gls{alptwo} & Alarmanlagen & Meldungen je Alarmanlage \\
        \hline
        Phase~2.1  & 3            & 10          & 3            & 100          & 100                      \\
        Phase~2.2  & 3            & \textbf{30} & 3            & 100          & 100                      \\
        Phase~2.3  & 3            & \textbf{60} & 3            & 100          & 100                      \\
        Phase~2.4  & 3            & 10          & \textbf{6}   & 100          & 100                      \\
        Phase~2.5  & 3            & 10          & \textbf{9}   & 100          & 100                      \\
        \hline
    \end{tabular}
\end{table}

Die folgenden Abbildungen zeigen den jeweiligen Ende-zu-Ende-Latenzverlauf; Tabelle~\ref{tab:amqp_exp_phase_2_latency} am Ende des Abschnitts fasst die Mittelwerte aller fünf Experimente zusammen.

\subsubsection*{Baseline (Phase~2.1)}
Abbildung~\ref{fig:amqp_exp_phase_2_1_latency} zeigt den Referenzverlauf bei zehn Queues und drei \gls{alptwo}-Instanzen. Dieser Verlauf dient als Vergleichsbasis für beide folgenden Skalierungsblöcke.

\begin{figure}[H]
    \centering
    \includegraphics[width=1\textwidth]{./img/exp/amqp/2/1/latency-line.pdf}
    \caption{Ende-zu-Ende-Latenzverlauf der Baseline (Phase~2.1).}
    \label{fig:amqp_exp_phase_2_1_latency}
\end{figure}

\subsubsection*{Queue-Skalierung (Phase~2.2 und Phase~2.3)}
Wird ausgehend von der Baseline ausschließlich die Queue-Anzahl auf 30 erhöht, sinkt die Latenz bereits deutlich (Abbildung~\ref{fig:amqp_exp_phase_2_2_latency}). Bei einer weiteren Erhöhung auf 60 Queues setzt sich dieser Trend fort (Abbildung~\ref{fig:amqp_exp_phase_2_3_latency}).

\begin{figure}[H]
    \centering
    \includegraphics[width=1\textwidth]{./img/exp/amqp/2/2/latency-line.pdf}
    \caption{Ende-zu-Ende-Latenzverlauf bei Queue-Skalierung auf 30 Queues (Phase~2.2).}
    \label{fig:amqp_exp_phase_2_2_latency}
\end{figure}

\begin{figure}[H]
    \centering
    \includegraphics[width=1\textwidth]{./img/exp/amqp/2/3/latency-line.pdf}
    \caption{Ende-zu-Ende-Latenzverlauf bei Queue-Skalierung auf 60 Queues (Phase~2.3).}
    \label{fig:amqp_exp_phase_2_3_latency}
\end{figure}

\subsubsection*{\gls{alptwo}-Skalierung (Phase~2.4 und Phase~2.5)}
Wird stattdessen bei unveränderter Queue-Anzahl die \gls{alptwo}-Kapazität auf sechs (Abbildung~\ref{fig:amqp_exp_phase_2_4_latency}) beziehungsweise neun Instanzen (Abbildung~\ref{fig:amqp_exp_phase_2_5_latency}) erhöht, bleibt der Latenzverlauf gegenüber der Baseline nahezu unverändert.

\begin{figure}[H]
    \centering
    \includegraphics[width=1\textwidth]{./img/exp/amqp/2/4/latency-line.pdf}
    \caption{Ende-zu-Ende-Latenzverlauf bei \acrshort{alptwo}-Skalierung auf 6 Instanzen (Phase~2.4).}
    \label{fig:amqp_exp_phase_2_4_latency}
\end{figure}

\begin{figure}[H]
    \centering
    \includegraphics[width=1\textwidth]{./img/exp/amqp/2/5/latency-line.pdf}
    \caption{Ende-zu-Ende-Latenzverlauf bei \acrshort{alptwo}-Skalierung auf 9 Instanzen (Phase~2.5).}
    \label{fig:amqp_exp_phase_2_5_latency}
\end{figure}

Tabelle~\ref{tab:amqp_exp_phase_2_latency} stellt die mittleren Ende-zu-Ende-Latenzen aller fünf Experimente gegenüber. Der Kontrast zwischen beiden Skalierungsblöcken wird hier besonders deutlich: Während die Latenz mit steigender Queue-Anzahl von 17\,654\,ms (10 Queues) auf 218\,ms (60 Queues) sinkt, verändert sich die Latenz bei zusätzlichen \gls{alptwo}-Instanzen bei gleichbleibender Queue-Anzahl kaum.

\begin{table}[H]
    \centering
    \caption{Zusammenfassung der Latenzmessungen in Phase~2 in Millisekunden.}
    \label{tab:amqp_exp_phase_2_latency}
    \begin{tabular}{lllll}
        \hline
        Experiment & Alarmmeldungen & Queues      & \gls{alptwo} & Mittelwert  \\
        \hline
        Phase~2.1  & 100            & 10          & 3            & 17\,654,122 \\
        Phase~2.2  & 100            & \textbf{30} & 3            & 2\,454,312  \\
        Phase~2.3  & 100            & \textbf{60} & 3            & 218,131     \\
        Phase~2.4  & 100            & 10          & \textbf{6}   & 17\,347,347 \\
        Phase~2.5  & 100            & 10          & \textbf{9}   & 17\,625,434 \\
        \hline
    \end{tabular}
\end{table}

\subsection{Ergebnisse und Interpretation}

Zusammengenommen bestätigen beide Phasen die in Abschnitt~\ref{sec:requirements} formulierten Anforderungen an \gls{alps}. Phase~1 zeigt, dass die Delivery-Garantien unter den untersuchten Fehlerbedingungen vollständig eingehalten werden: Die Validator-Auswertung weist weder Verletzungen der \textit{Exactly-Once-Delivery} noch der \textit{In-Order-Delivery} noch Lücken in der Vollständigkeit aus.

Für die Skalierbarkeit liefert Phase~2 eine klare Erklärung für den beobachteten Unterschied zwischen Queue- und Instanzskalierung. Durch \texttt{prefetch\_count = 1} und \texttt{x-single-active-consumer} ist pro Queue jeweils nur ein aktiver \gls{alptwo}-Verarbeitungspfad wirksam. Dadurch kann eine langsame Nachricht alle nachfolgenden Nachrichten derselben Queue verzögern, ein Effekt, der als \textit{Head-of-Line-Blocking}\footnote{Head-of-Line-Blocking bezeichnet den Effekt, dass eine frühe, langsame Nachricht nachfolgende Nachrichten derselben Queue blockiert~\citep{kurose2017computer}.} bezeichnet wird. Zusätzliche \gls{alptwo}-Instanzen ändern daran wenig, solange die Anzahl der Queues konstant bleibt, weil weiterhin nur eine begrenzte Anzahl an Queues parallel abgearbeitet werden kann. Eine höhere Queue-Anzahl hingegen reduziert dieses Blocking direkt, da mehr unabhängige Verarbeitungspfade zur Verfügung stehen. Dies erklärt die in Tabelle~\ref{tab:amqp_exp_phase_2_latency} beobachtete starke Latenzreduktion.

Aus betrieblicher Sicht bedeutet dies, dass Skalierung in \gls{alps} überwiegend über brokerseitige Queue-Parallelisierung erfolgt und weniger über zusätzliche \gls{alptwo}-Services. Im Vergleich zum \gls{mqtt}-basierten \gls{lps}, das primär über zusätzliche Service-Instanzen skalierte, spricht dies für einen geringeren operativen Skalierungsaufwand bei wachsender Last: Statt neue Services zu deployen und zu überwachen, genügt eine Anpassung der Queue-Konfiguration am Broker. Wie sich diese Unterschiede zwischen \gls{lps} und \gls{alps} insgesamt einordnen lassen, wird im folgenden Kapitel im Detail diskutiert.

\chapter{Vergleich und Diskussion}
\label{ch:vergleich}

Dieses Kapitel bewertet die in dieser Arbeit entwickelte \gls{alps}-Architektur. Im ersten Teil wird \gls{alps} dem ursprünglichen \gls{lps} gegenübergestellt und anhand der experimentellen Ergebnisse bewertet. Anschließend wird diskutiert, inwiefern \gls{alps} die abgeleiteten Anforderungen erfüllt und wo die Grenzen der vorliegenden Evaluation liegen.

\section{Vergleich: \acrshort{lps} vs. \acrshort{alps}}

Beide Architekturen verfolgen das primäre Ziel, synchrone Request-Response-Ketten eines verteilten Legacy-Systems aufzubrechen. Beide Ansätze erreichen diese asynchrone Entkopplung erfolgreich, ohne dass der bestehende Quellcode der Legacy-Services angepasst werden muss. Der entscheidende Unterschied liegt jedoch in der Korrektheit unter Fehlerbedingungen und im Konzept der Skalierbarkeit.

\subsection{Umgang mit Fehlern und Delivery-Garantien}
Die Evaluation hat deutlich gezeigt, dass das ursprüngliche \gls{lps} zwar im fehlerfreien Betrieb funktioniert, unter Fehlerbedingungen (Netzwerkfehler, Serverausfälle) jedoch fundamentale Delivery-Garantien verletzt. Im \gls{lps}-Experiment traten bei $5\,000$ gesendeten Alarmmeldungen $17\,504$ Duplikate ($77{,}8\,\%$) und $3\,035$ Reihenfolgeverletzungen ($13{,}5\,\%$) auf. Die Ursache hierfür ist, dass das \gls{lps} sich ausschließlich auf brokerseitige Retry-Mechanismen verlässt. Ohne einen echten anwendungsseitigen Zustand können fachliche Duplikate nicht erkannt werden, und die Reihenfolge gerät durch parallele Retries durcheinander.

\gls{alps} löst dieses Problem durch eine funktionale Trennung:

Die \textit{In-Order-Delivery} wird vollständig vom Broker übernommen. Durch deterministisches Routing (\texttt{x-consistent-hash}) landen alle Nachrichten derselben Alarmanlage in derselben Queue. Die Kombination aus \texttt{x-single-active-consumer} und \texttt{prefetch\_count = 1}garantiert, dass bei Fehlern (\gls{amqp}-NACK) eine Nachricht so lange blockiert, bis sie erfolgreich verarbeitet wurde. Spätere Nachrichten können sie nicht überholen.

Die \textit{Exactly-Once-Delivery} wird durch die nachgelagerte Sidecar-Komponente \gls{alpthree} abgesichert. Durch die verteilte Mnesia-Datenbank entsteht ein fachliches Gedächtnis unmittelbar vor dem Legacy-Service (\gls{srv}~3), das Duplikate zuverlässig herausfiltert.

Die quantitativen Ergebnisse belegen diesen Erfolg: Unter exakt denselben künstlich induzierten Fehlerbedingungen, unter denen das \gls{lps} massiv versagte, wies \gls{alps} in der Validierung $0\,\%$ Duplikate und $0\,\%$ Reihenfolgeverletzungen auf.

\subsection{Unterschiede in der Skalierbarkeit}
Ein weiterer wesentlicher Unterschied zeigt sich im operativen Betrieb unter Last. Das \gls{lps} skalierte in den Experimenten primär über das Hinzufügen weiterer Service-Instanzen, was den Betriebs- und Wartungsaufwand erhöht.

\gls{alps} hingegen verlagert die Parallelisierung auf die Broker-Ebene. Die Latenzmessungen in Phase~2 (vgl. Abschnitt~\ref{sec:eval-skalierbarkeit}) haben gezeigt, dass eine Erhöhung der Queue-Anzahl (von 10 auf 60 Queues) bei gleichbleibender Consumer-Anzahl zu einer signifikanten Reduktion der Ende-zu-Ende-Latenz führt. Zusätzliche \gls{alptwo}-Instanzen verbessern die Durchlaufzeiten zwar ebenfalls, der primäre Hebel für die Leistungssteigerung liegt im \gls{alps}-Design jedoch in der effizienten Partitionierung durch Queues. Dies bedeutet in der Praxis einen geringeren administrativen Aufwand, da zur Lastbewältigung weniger neue Microservices deployt und überwacht werden müssen.

\section{Erfüllung der Anforderungen}

In Abschnitt~\ref{sec:requirements}  wurden acht konkrete Anforderungen (A1-A8) an das neue System gestellt. Das entwickelte \gls{alps} erfüllt diese Vorgaben durch folgende Architekturentscheidungen:

\begin{itemize}
    \item \textbf{A1 (Asynchrone Entkopplung):} Die asynchrone Entkopplung wird dadurch erreicht, dass die Komponente \gls{alpone} eingehende HTTP-Requests sofort bestätigt, sobald das Event erfolgreich an den Broker übergeben wurde. Die vorgelagerten Systeme müssen nicht mehr auf die vollständige Verarbeitung der gesamten Kette warten.

    \item \textbf{A2 (\textit{In-Order-Delivery} pro Alarmanlage):} Die strikte Reihenfolge pro Alarmanlage wird durch \texttt{x-single-active-consumer} in Kombination mit \texttt{prefetch\_count = 1} sichergestellt. So wird garantiert, dass bei temporären Fehlern eine nachfolgende Nachricht derselben Anlage die fehlerhafte nicht überholen kann.

    \item \textbf{A3 (Fachliche Duplikate robust behandeln):} \gls{alps} löst das Duplikat-Problem durch das zusätzliche \gls{alpthree} vor \gls{srv}~3. Diese Schicht nutzt eine verteilte Mnesia-Datenbank, um wiederholte Auslieferungen anhand der Kombination aus \texttt{device\_id} und \texttt{triggertime} zu erkennen und zu blockieren, bevor sie den Legacy-Service erreichen.

    \item \textbf{A4 (Fehlerfälle ohne Nachrichtenverlust behandeln):} Tritt bei der nachgelagerten Verarbeitung ein Fehler auf (z.\,B. Netzwerkabbruch oder Serverfehler), sendet \gls{alptwo} ein \gls{amqp}-NACK mit dem Parameter \texttt{requeue=true}. Die Nachricht geht dadurch nicht verloren, sondern wird vom Broker sicher zur erneuten Verarbeitung bereitgestellt.

    \item \textbf{A5 (Kontrollierte Parallelisierung ermöglichen):} Das System nutzt einen Exchange vom Typ \texttt{x-consistent-hash}, der Nachrichten anhand der \texttt{device\_id} aufteilt. Dadurch können Nachrichten unterschiedlicher Alarmanlagen über mehrere Queues parallel verarbeitet werden, während die Logik aus A2 gleichzeitig die Reihenfolge für jede einzelne Alarmanlage schützt.

    \item \textbf{A6 (Betriebsaufwand bei Skalierung begrenzen):} Die Experimente zur Skalierbarkeit zeigen, dass Laststeigerungen im \gls{alps} effizient durch eine reine Erhöhung der Queue-Anzahl abgefedert werden können. Dies reduziert den operativen Aufwand, da zur Skalierung keine neuen Service-Instanzen (\gls{alptwo}) deployt und überwacht werden müssen.

    \item \textbf{A7 (Vendor-Lock-in reduzieren):} Im Gegensatz zum \gls{lps}, das für seine Fehlerbehandlung auf ein sehr spezifisches HiveMQ-Plugin angewiesen ist, nutzt \gls{alps} standardnahe \gls{amqp}-Protokollfunktionen. Die komplexe Logik zur Deduplizierung wurde zudem sauber in eine eigene Komponente (\gls{alpthree}) ausgelagert, was einen zukünftigen Austausch des Message Brokers deutlich vereinfacht.

    \item \textbf{A8 (Transparenz zum Legacy-System bewahren):} Da \gls{alpone} nach außen hin genau dieselbe HTTP-Schnittstelle bereitstellt wie der ursprüngliche Legacy-Service, kann \gls{alps} nahtlos zwischen die bestehenden Systeme geschaltet werden. Das Legacy-System merkt von der neuen Zwischenschicht nichts.
\end{itemize}

\section{Limitationen der Experimente}

Obwohl die Experimente die funktionale Korrektheit und Skalierbarkeit von \gls{alps} belegen, unterliegt die Evaluation gewissen Einschränkungen, die bei einer Übertragung in eine reale Produktionsumgebung berücksichtigt werden müssen:

\begin{itemize}
    \item \textbf{Synthetische Fehler- und Lastszenarien:} Die Experimente decken definierte Fehlerraten (z. B. $0,1\,\%$ \gls{tcp}-Drops und Serverfehler) und spezifische Lastprofile ab. Extreme Ausnahmezustände, wie beispielsweise ein plötzlicher Ausfall kompletter Rechenzentren unter Maximallast oder stark schwankende Netzwerklatenzen, wurden nicht simuliert.

    \item \textbf{Abhängigkeit von Systemparametern:} Die ermittelten Latenzzeiten und die Skalierungseffekte sind eng an die gewählten Parameter und Hardware-Spezifikationen gebunden. In einer realen Cloud- oder On-Premise-Umgebung werden die absoluten Latenzwerte abweichen.

    \item \textbf{Auflösung von Split-Brain-Szenarien:} Mnesia wurde so konfiguriert, dass es isolierte Datenbank-Knotenpunkte nach einer Netzwerkpartitionierung (\textit{Split-Brain}) mittels eines Append-Only-Verfahrens wieder zusammenführt. Während dies theoretisch sicherstellt, dass keine Daten gelöscht werden, wurde das tatsächliche Verhalten des Mnesia-Clusters unter andauernden, instabilen Netzwerkabbrüchen im Rahmen dieser Experimente nicht abschließend getestet.
\end{itemize}

Zusammenfassend zeigt die Evaluation jedoch eindeutig, dass \gls{alps} die in der Vorarbeit identifizierten Limitierungen erfolgreich behebt. Es bietet eine architektonisch saubere, skalierbare und vor allem fehlerresistente Lösung, um geschäftskritische Legacy-Systeme sicher zu modernisieren.

\chapter{Fazit}
\label{ch:fazit}

Die vorliegende Masterarbeit befasste sich mit der strukturellen Optimierung verteilter Legacy-Systeme, die aufgrund historisch gewachsener Design Smells zunehmend schwer wartbar und fehleranfällig geworden sind. Das zentrale Ziel war es, ein Architektur-Reengineering durchzuführen, das synchrone Kommunikationsketten aufbricht und moderne Anforderungen an Skalierbarkeit sowie Korrektheit erfüllt, ohne dass der Quellcode der zugrunde liegenden Legacy-Services modifiziert werden muss.

Ausgangspunkt der Untersuchung war das bestehende \gls{lps}, ein \gls{mqtt}-basierter Architekturansatz. Die systematische experimentelle Analyse in dieser Arbeit deckte auf, dass das \gls{lps} im störungsfreien Betrieb zwar funktionale Vorteile bietet, unter Fehlerbedingungen jedoch zentrale Delivery-Garantien des Systems verletzt. Insbesondere fehlten dem Ansatz ausreichende Mechanismen, um fachliche Duplikate (\textit{Exactly-Once-Delivery}) zu verhindern und die strikte Verarbeitungsreihenfolge (\textit{In-Order-Delivery}) unter Last sowie bei Netzwerkausfällen verlässlich zu wahren. Zudem erwies sich die Skalierung primär über zusätzliche Service-Instanzen als operativ aufwändig.

Um diese Limitierungen zu überwinden, wurde in dieser Arbeit das \gls{alps} konzipiert und implementiert. Durch den Wechsel auf ein \gls{amqp}-basiertes Messaging-Protokoll und den gezielten Einsatz von Broker-Mechanismen wie \texttt{x-consistent-hash}, \texttt{x-single-active-consumer} sowie einer strikten Flusskontrolle (\texttt{prefetch\_count = 1}) konnte die \textit{In-Order-Delivery} deterministisch pro Alarmanlage sichergestellt werden. Ergänzt wurde die Architektur durch eine vorgelagerte Sidecar-Komponente (\gls{alpthree}) auf Basis der verteilten Erlang-Datenbank Mnesia. Diese Deduplizierungsschicht fungiert als fachliches Gedächtnis und filtert Duplikate robust heraus, bevor sie den eigentlichen Legacy-Service erreichen.

Die anschließende experimentelle Evaluation bestätigte den Erfolg dieses Architektur-Refactorings. Unter exakt jenen künstlich induzierten Fehlerbedingungen, die beim ursprünglichen \gls{lps} zu massiven Fehlfunktionen führten (über 77\,\% Duplikate und 13\,\% Reihenfolgeverletzungen), wies \gls{alps} eine Fehlerquote von $0\,\%$ auf. Beide Delivery-Garantien wurden ausnahmslos eingehalten. Darüber hinaus zeigte die Skalierbarkeitsanalyse, dass \gls{alps} Laststeigerungen effizient durch reine Queue-Partitionierung abfedern kann. Dies reduziert den Bedarf an zusätzlichen \gls{alptwo}-Instanzen und minimiert den Betriebsaufwand unter Hochlast.

\chapter{Future Work}
\label{ch:future-work}

Obwohl die Ergebnisse dieser Masterarbeit die funktionale Korrektheit und Skalierbarkeit von \gls{alps} belegen, eröffnen die getroffenen Architekturentscheidungen und die Limitierungen der experimentellen Evaluation konkrete Anknüpfungspunkte für zukünftige Forschungsarbeiten. Diese lassen sich in drei zentrale Bereiche unterteilen:

\section{Evaluation unter realen Cloud- und On-Premise-Bedingungen}
Die in Kapitel~\ref{ch:eval} durchgeführten Experimente basierten auf lokal orchestrierten Docker-Compose-Umgebungen. Obwohl diese synthetischen Setups eine reproduzierbare Fehlerinjektion ermöglichen, spiegeln sie die Komplexität physisch verteilter Infrastrukturen nicht vollständig wider. Zukünftige Arbeiten sollten \gls{alps} in realen, geographisch verteilten Cloud- oder On-Premise-Umgebungen ausrollen. Dabei gilt es, den Einfluss von realen Netzwerklatenzen, Bandbreitenbeschränkungen und variierenden Hardware-Ressourcen  zu untersuchen. Solche Feldtests sind entscheidend, um die produktive Einsatzreife für industrielle Hochverfügbarkeitsumgebungen abschließend zu validieren.

\section{Validierung des Mnesia-Clusterings unter Netzwerkpartitionierung}
Die Wahrung der \textit{Exactly-Once-Delivery} stützt sich maßgeblich auf die Sidecar-Komponente \gls{alpthree} und die zugrunde liegende verteilte Mnesia-Datenbank. Wie in Abschnitt~\ref{subsec:grundlagen_mnesia} diskutiert, wurde Mnesia so konfiguriert, dass isolierte Knoten nach einem Split-Brain-Szenario mittels eines Append-Only-Verfahrens wieder zusammengeführt werden. Das tatsächliche Verhalten dieses Mnesia-Clusters unter andauernden, instabilen Netzwerkabbrüchen zwischen den Cluster-Knoten wurde in den bisherigen Experimenten jedoch nicht empirisch validiert. Zukünftige Untersuchungen müssen gezielte Netzwerkpartitionierungen auf Datenbankebene induzieren, um die Zuverlässigkeit des \textit{unsplit}-Mechanismus zu prüfen.

\section{Erweiterte Simulationen mit \acrshort{rast}}
\label{sec:future-work-rast}

Ein weiterer essenzieller Schritt zur tiefergehenden Validierung der \gls{alps}-Architektur ist die Nutzung des \gls{rast}~\citep{rast}. Dieses Toolset wurde gezielt für die Performance-Vorhersage von zeitkritischen Legacy-Systemen entwickelt, die im produktiven Einsatz stehen und nicht durch aufwändige Application Performance Monitoring Tools beeinträchtigt werden dürfen~\citep{rast}. Stattdessen extrahiert \gls{rast} historische Leistungsdaten aus einfachen Request-Logs und nutzt maschinelle Lernverfahren, um Vorhersagemodelle für die Verarbeitungszeiten zu trainieren~\citep{rast}. Diese Modelle werden anschließend in einer Simulator-Komponente ausgeführt, die mittels konfigurierbarer Lastintensitätsprofile das Antwortzeitverhalten des Systems unter realitätsnahen Bedingungen misst~\citep{rast}.

Für die zukünftige Forschung bietet die formale Integration von \gls{alps} in das \gls{rast}-Toolset den entscheidenden Vorteil, das architektonische Leistungsverhalten unter hochkomplexen Lastszenarien prädiktiv bewerten zu können, ohne die reale Produktionsumgebung zu gefährden. Während sich die bisherige Evaluation auf deterministische Fehlerinjektionen in einer isolierten Docker-Umgebung beschränkte, ließe sich mit \gls{rast} simulieren, wie sich die neu eingeführte \gls{amqp}-Schicht und die Mnesia-basierte Deduplizierung in ressourcenteilenden Umgebungen verhalten.

\appendix

\chapter{Anhang}
\label{ch:anhang}

\section{Reproduzierbarkeit der Experimente}
\label{sec:appendix-reproducibility}

Die in Abschnitt~\ref{sec:mqtt-proxy-experimente} und Kapitel~\ref{ch:eval} beschriebenen Experimente sind reproduzierbar, indem die benötigten Docker-Images geladen, eigene Konfigurationsdateien erstellt und die Simulationen anschließend gestartet werden.

\subsection{Voraussetzungen}

\begin{itemize}
    \item Installiertes Docker~\citep{docker_2023}.
    \item Installiertes Docker Compose~\citep{docker_compose}.
\end{itemize}

\subsection{Docker-Images laden}

\begin{verbatim}
docker pull rabbitmq:4-management
docker pull jadkhaddad/lps-rs
docker pull jadkhaddad/lps-gleam
docker pull jadkhaddad/lps-scripts
docker pull jadkhaddad/lps-mqtt-python
docker pull jadkhaddad/lps-mqtt-hivemq
\end{verbatim}

\subsection{Bedeutung der Docker-Images}

\begin{itemize}
    \item \texttt{rabbitmq:4-management}: RabbitMQ-Broker für das \gls{amqp}-Szenario.
    \item \texttt{jadkhaddad/lps-rs}: Enthält die Simulation-Runner sowie die Kern-Komponenten (u.\,a. Alarmanlagen-Simulator, \gls{alpone}, \gls{lptwo}, Server-Simulator, Validator, \gls{tcp}-Proxy).
    \item \texttt{jadkhaddad/lps-gleam}: Enthält die Deduplizierungs-Komponente (\gls{alpthree}) für \gls{alps}.
    \item \texttt{jadkhaddad/lps-scripts}: Enthält Auswertungsskripte zur Plot-Erzeugung.
    \item \texttt{jadkhaddad/lps-mqtt-python}: Enthält die \gls{mqtt}-bezogenen Proxy-Skripte (u.\,a. \gls{srv}~1, \gls{srv}~2, \gls{lpone}, \gls{lptwo}).
    \item \texttt{jadkhaddad/lps-mqtt-hivemq}: Enthält den HiveMQ-Broker mit den verwendeten Plugins für das \gls{mqtt}-Szenario.
\end{itemize}

\subsection{Docker-Netzwerk erstellen}

Vor dem Start wird ein Docker-Netzwerk angelegt, über das die Komponenten mit einander kommunizieren.

\begin{verbatim}
docker network create lps-simulation-net
\end{verbatim}

\subsection{Konfigurationsdateien anlegen}

Vor dem Start werden die benötigten Simulationskonfigurationen als \gls{yaml}-Dateien angelegt (z.\,B. eine \gls{amqp}-Konfiguration und eine \gls{mqtt}-Konfiguration).
Diese Dateien enthalten die jeweiligen Parameter für Last, Fehlerprofile und Laufzeit der Experimente.

\subsubsection*{Beispiel: \gls{amqp}-Konfiguration}

\begin{verbatim}
simulation_dir: "simulation"
simulation_docker_network: "lps-simulation-net"

lps_rust_image: "jadkhaddad/lps-rs"
lps_scripts_image: "jadkhaddad/lps-scripts"
lps_gleam_image: "jadkhaddad/lps-gleam"

devices:
  devices_count: 100
  max_event_interval: 500ms
  retry_duration: 500ms
  max_events: 100

server:
  failure_rate: 0.1
  max_work_duration: 50ms
  max_response_duration: 50ms
  deduplication: true
  server_count: 3

amqp:
  queues_count: 10
  producers:
    producers_count: 3
  consumers:
    consumers_count: 3
    reconnect_duration: 500ms
    amqp_proxy:
      failure_rate: 0.1
      max_up_duration: 10s
\end{verbatim}

\subsubsection*{Beispiel: \gls{mqtt}-Konfiguration}

\begin{verbatim}
simulation_dir: "simulation-mqtt"
simulation_docker_network: "lps-simulation-net"

lps_rust_image: "jadkhaddad/lps-rs"
lps_scripts_image: "jadkhaddad/lps-scripts"
lps_gleam_image: "jadkhaddad/lps-gleam"
lps_mqtt_python_image: "jadkhaddad/lps-mqtt-python"
lps_mqtt_hivemq_image: "jadkhaddad/lps-mqtt-hivemq"

devices:
  devices_count: 100
  max_event_interval: 500ms
  retry_duration: 500ms
  max_events: 100

server:
  failure_rate: 0.1
  max_work_duration: 50ms
  max_response_duration: 50ms
  server_count: 3

mqtt:
  producers:
    producers_count: 3
  consumers:
    consumers_count: 3
    mqtt_proxy:
      failure_rate: 0.1
      max_up_duration: 10s
\end{verbatim}

\subsubsection*{Bedeutung der Parameter}

\paragraph{Allgemein}
\begin{itemize}
    \item \texttt{simulation\_dir}: Zielverzeichnis der erzeugten Simulationsartefakte.
    \item \texttt{simulation\_docker\_network}: Externes Docker-Netzwerk, in dem die Services gestartet werden.
    \item \texttt{lps\_rust\_image}: Image-Name für die Simulation-Runner und Kernkomponenten.
    \item \texttt{lps\_scripts\_image}: Image-Name für die Auswertungsskripte.
    \item \texttt{lps\_gleam\_image}: Image-Name für die Deduplizierungs-Komponente (\gls{alpthree}).
    \item \texttt{lps\_mqtt\_python\_image}: Image-Name für die Python-basierten \gls{mqtt}-Proxy-Komponenten.
    \item \texttt{lps\_mqtt\_hivemq\_image}: Image-Name für den HiveMQ-Broker im \gls{mqtt}-Szenario.
\end{itemize}

\paragraph{Block \texttt{devices}}
\begin{itemize}
    \item \texttt{devices\_count}: Anzahl simulierter Alarmanlagen.
    \item \texttt{max\_event\_interval}: Maximales Intervall zwischen zwei erzeugten Meldungen pro Alarmanlage.
    \item \texttt{retry\_duration}: Wartezeit bis zum erneuten Sendeversuch bei fehlgeschlagener Übertragung.
    \item \texttt{max\_events}: Maximale Anzahl erzeugter Meldungen pro Alarmanlage.
\end{itemize}

\paragraph{Block \texttt{server}}
\begin{itemize}
    \item \texttt{failure\_rate}: Wahrscheinlichkeit für Fehlerantworten des Server-Simulators.
    \item \texttt{max\_work\_duration}: Maximale Bearbeitungszeit einer Anfrage im Server.
    \item \texttt{max\_response\_duration}: Maximale zusätzliche Antwortverzögerung.
    \item \texttt{server\_count}: Anzahl gestarteter Server-Instanzen.
    \item \texttt{deduplication} (optional): Aktiviert die zusätzliche Deduplizierungs-Komponente (\gls{alps}-Szenario).
\end{itemize}

\paragraph{Block \texttt{amqp}}
\begin{itemize}
    \item \texttt{queues\_count}: Anzahl der Broker-Queues.
    \item \texttt{producers.producers\_count}: Anzahl \gls{amqp}-Producer-Instanzen (\gls{alpone}).
    \item \texttt{consumers.consumers\_count}: Anzahl \gls{amqp}-Consumer-Instanzen (\gls{alptwo}).
    \item \texttt{consumers.reconnect\_duration}: Wartezeit vor Reconnect eines Consumers nach Verbindungsabbruch.
    \item \texttt{consumers.amqp\_proxy} (optional): Fehlerproxy zwischen Consumer und Broker.
          \begin{itemize}
              \item \texttt{failure\_rate}: Wahrscheinlichkeit eines künstlichen Verbindungsfehlers.
              \item \texttt{max\_up\_duration}: Maximale Dauer einer Verbindung vor erzwungener Trennung.
          \end{itemize}
\end{itemize}

\paragraph{Block \texttt{mqtt}}
\begin{itemize}
    \item \texttt{producers.producers\_count}: Anzahl \gls{mqtt}-Producer-Instanzen (\gls{lpone}).
    \item \texttt{consumers.consumers\_count}: Anzahl \gls{mqtt}-Consumer-Instanzen (\gls{lptwo}).
    \item \texttt{consumers.mqtt\_proxy} (optional): Fehlerproxy zwischen \gls{mqtt}-Consumer-Pfad und Broker.
          \begin{itemize}
              \item \texttt{failure\_rate}: Wahrscheinlichkeit eines künstlichen Verbindungsfehlers.
              \item \texttt{max\_up\_duration}: Maximale Dauer einer Verbindung vor erzwungener Trennung.
          \end{itemize}
\end{itemize}

\paragraph{Dauerangaben}
Zeitwerte werden als Dauerangaben angegeben (z.\,B. \texttt{500ms}, \texttt{10s}).

\subsection{Simulationen starten}

\subsubsection*{\gls{amqp}}

\begin{verbatim}
docker run --rm -it \
  -v /var/run/docker.sock:/var/run/docker.sock \
  -v "$(pwd):/app" \
  jadkhaddad/lps-rs \
  simulation-runner \
  --config-file <amqp-config.yaml> \
  --root-abs-path $(pwd)
\end{verbatim}

\subsubsection*{\gls{mqtt}}

\begin{verbatim}
docker run --rm -it \
  -v /var/run/docker.sock:/var/run/docker.sock \
  -v "$(pwd):/app" \
  jadkhaddad/lps-rs \
  mqtt-simulation-runner \
  --config-file <mqtt-config.yaml> \
  --root-abs-path $(pwd)
\end{verbatim}

\subsection{Ergebnisse}

Nach Abschluss eines Laufs befinden sich die erzeugten Ergebnisse unter\newline \texttt{\$(pwd)/mount}.

\section{Quellcode}
\label{s:source-code}

Dieses Projekt ist ein Open-Source-Projekt~\citep{source} und unter der \textit{MIT License}~\citep{MIT} lizenziert.

\begin{figure}[h]
    \centering
    \includegraphics[width=0.4\textwidth]{./img/qrcode.png}
    \caption{Quellcode}
    \label{fig:sc}
\end{figure}